\documentclass[11pt, a4paper]{article}

% --- Fonts and engine ---
\usepackage{fontspec}
\setmainfont{Helvetica Neue}
\setmonofont{Menlo}
\usepackage{unicode-math}
\setmathfont{STIX Two Math}

% --- Layout ---
\usepackage[margin=2cm, top=2cm, bottom=2.5cm]{geometry}
\usepackage{microtype}
\usepackage{parskip}
\setlength{\parskip}{0.45em}
\setlength{\parindent}{0pt}

% --- Typography and links ---
\usepackage[colorlinks=true, linkcolor=NavyBlue, urlcolor=NavyBlue, citecolor=NavyBlue]{hyperref}
\usepackage{xcolor}
\usepackage[dvipsnames]{xcolor}

% --- Math and tables ---
\usepackage{amsmath}
\usepackage{booktabs}
\usepackage{array}
\usepackage{tabularx}
\newcommand{\thead}[1]{\textbf{#1}}

% --- Lists ---
\usepackage{enumitem}
\setlist[itemize]{leftmargin=1.4em, itemsep=0.1em, topsep=0.2em}
\setlist[enumerate]{leftmargin=1.5em, itemsep=0.1em, topsep=0.2em}

% --- Boxed callouts ---
\usepackage[most]{tcolorbox}
\tcbuselibrary{breakable, skins, theorems}

\definecolor{defBg}{HTML}{EAF2FB}
\definecolor{defBd}{HTML}{1F4E79}
\definecolor{tipBg}{HTML}{E8F5E9}
\definecolor{tipBd}{HTML}{2E7D32}
\definecolor{trapBg}{HTML}{FCE4EC}
\definecolor{trapBd}{HTML}{AD1457}
\definecolor{pitBg}{HTML}{FFF8E1}
\definecolor{pitBd}{HTML}{B27600}
\definecolor{noteBg}{HTML}{F3E5F5}
\definecolor{noteBd}{HTML}{6A1B9A}
\definecolor{tldrBg}{HTML}{E1F5FE}
\definecolor{tldrBd}{HTML}{0277BD}
\definecolor{exBg}{HTML}{F4F4F4}
\definecolor{exBd}{HTML}{555555}

\newtcolorbox{defbox}[1][]{enhanced, breakable, colback=defBg, colframe=defBd, arc=2pt, boxrule=0.6pt, left=8pt, right=8pt, top=5pt, bottom=5pt, fonttitle=\bfseries, title={Definition: #1}}
\newtcolorbox{exbox}[1][]{enhanced, breakable, colback=exBg, colframe=exBd, arc=2pt, boxrule=0.6pt, left=8pt, right=8pt, top=5pt, bottom=5pt, fonttitle=\bfseries, title={Worked example: #1}}
\newtcolorbox{setupbox}[1][]{enhanced, breakable, colback=defBg, colframe=defBd, arc=2pt, boxrule=0.6pt, left=8pt, right=8pt, top=5pt, bottom=5pt, fonttitle=\bfseries, title={Setup: #1}}
\newtcolorbox{exambox}[1][]{enhanced, breakable, colback=tipBg, colframe=tipBd, arc=2pt, boxrule=0.6pt, left=8pt, right=8pt, top=5pt, bottom=5pt, fonttitle=\bfseries, title={Exam-mode answer #1}}
\newtcolorbox{trapbox}[1][]{enhanced, breakable, colback=trapBg, colframe=trapBd, arc=2pt, boxrule=0.6pt, left=8pt, right=8pt, top=5pt, bottom=5pt, fonttitle=\bfseries, title={MCQ trap #1}}
\newtcolorbox{pitfallbox}[1][]{enhanced, breakable, colback=pitBg, colframe=pitBd, arc=2pt, boxrule=0.6pt, left=8pt, right=8pt, top=5pt, bottom=5pt, fonttitle=\bfseries, title={Pitfall #1}}
\newtcolorbox{notebox}[1][]{enhanced, breakable, colback=noteBg, colframe=noteBd, arc=2pt, boxrule=0.6pt, left=8pt, right=8pt, top=5pt, bottom=5pt, fonttitle=\bfseries, title={Lecturer aside #1}}
\newtcolorbox{tldrbox}[1][]{enhanced, breakable, colback=tldrBg, colframe=tldrBd, arc=2pt, boxrule=0.6pt, left=8pt, right=8pt, top=5pt, bottom=5pt, fonttitle=\bfseries, title={#1}}

\usepackage{titlesec}
\titleformat{\section}[block]{\Large\bfseries\color{defBd}}{\thesection.}{0.6em}{}
\titleformat{\subsection}[block]{\large\bfseries\color{defBd!80!black}}{\thesubsection}{0.5em}{}
\titleformat{\subsubsection}[block]{\normalsize\bfseries}{}{0pt}{}
\titlespacing*{\section}{0pt}{1.4em}{0.45em}
\titlespacing*{\subsection}{0pt}{0.9em}{0.3em}

\usepackage{fancyhdr}
\pagestyle{fancy}
\fancyhf{}
\fancyhead[L]{\small CM52054 — Combined Revision}
\fancyhead[R]{\small Weeks 1–4, 7–11, 19–23, 25–26, 29}
\fancyfoot[C]{\small\thepage}
\renewcommand{\headrulewidth}{0.3pt}

\providecommand{\blacksquare}{\rule{0.55em}{0.55em}}
\newcommand{\examflag}{\textcolor{tipBd}{\textbf{[EXAM]}}\,}
\newcommand{\trapflag}{\textcolor{trapBd}{\textbf{[TRAP]}}\,}
\newcommand{\kw}[1]{\textbf{#1}}

\title{\vspace{-1.2cm}{\Huge\bfseries Combined Revision Notes} \\[0.4em]{\large Trees (1–4) $\cdot$ Optimisation (7–9) $\cdot$ Bayesian \& Logistic (10–11) $\cdot$ Unsupervised (19) $\cdot$ Parametric (20) $\cdot$ Regularisation (21) $\cdot$ Graphical Models (22–23) $\cdot$ SVM (25–26) $\cdot$ Kernels (29)}}
\author{\large CM52054 Foundational Machine Learning \\[0.2em]\normalsize University of Bath}
\date{Condensed mark-scheme-style revision}

\begin{document}

\maketitle
\thispagestyle{empty}

\begin{tldrbox}[How to use]
The complete final revision document — Weeks 1–4, 7–11, 19–23, 25–26. Statements are at
\emph{mark-scheme keyword} density: every definition and setup is boxed, every examinable derivation
is included inline, pitfalls and MCQ traps are flagged. Non-examinable material is omitted; where a
\emph{proof} is non-examinable but its \emph{idea} is testable, only the idea is given. Revise
directly from this document — the one-line recall checklist at the end is the last-minute pass.
\end{tldrbox}

\begin{center}\large\bfseries\color{defBd} PART I — Foundations, Trees, Bias–Variance, Performance\end{center}

% =====================================================================
\section{Week 1 — the chassis}

Week~1 owns no past-paper question, but it sets up the vocabulary every later question silently
assumes — so treat it as the conceptual chassis the rest of the unit bolts onto.

The starting move is to distinguish three ways a computer can act intelligently. In
\kw{programming} a human writes the rule directly — every threshold and branch is hand-coded — so the
machine ``does exactly what you tell it and nothing else.'' In \kw{classical AI} the human specifies
the \emph{problem} and the computer \emph{searches} a fully built problem space for the best solution
(A$^\ast$ on a road graph). In \kw{ML} the computer gets neither the rule nor the full problem space:
it is handed input/output \emph{examples} and must \emph{induce} the rule itself. The relationship is
strictly nested: ML $\subset$ AI $\subset$ Programming — every ML system is doing AI, every AI system
is a program, but not every program learns.

\begin{itemize}
  \item \kw{Five-step pipeline}: (1) choose problem; (2) obtain data; (3) choose/design model;
  (4) fit by optimisation; (5) measure performance — then iterate. The loop matters: a poor
  performance measurement feeds straight back into re-choosing or re-fitting the model. This unit
  ignores step~2 and concentrates on steps~3--5.
  \item Model fitting \kw{is} optimisation: pick a parametric family $f_{\boldsymbol\theta}$, pick a
  loss $L(\boldsymbol\theta)$ scoring how badly it fits, then \emph{search} for
  $\boldsymbol\theta^\ast=\arg\min L$. This single schema — family, loss, search — recurs for every
  model in the unit; only the family and the optimiser change.
  \item \kw{Classification} = discrete label output (fish vs invertebrate, spam vs not); \kw{regression}
  = real-valued target (tomorrow's stock price). This is the \emph{type of output}; a separate axis is
  whether the output is a hard label or a probability.
  \item \kw{Train/test split}: fit on the training set, judge on an unseen test set. The reason this
  is non-negotiable is that several different rules can fit a small training set perfectly yet
  disagree on test data — perfect training accuracy therefore guarantees \emph{nothing} about
  generalisation. The lecturer's ``bear'' dramatises the sharpest version: deployed in the wild the
  model meets a mammal, an input its fish/invertebrate family simply cannot represent.
\end{itemize}

\begin{notebox}[an ML system is code plus a parameter blob]
The lecturer's slogan ``an ML algorithm outputs code'' is best read as: the \emph{architecture} (the
function that consumes a feature vector and returns a prediction) is fixed at design time, and
training only produces the \emph{parameters} that, plugged into that fixed function, pick out one
specific predictor. A trained model is therefore exactly two things — static code plus a parameter
blob — which is the working picture behind every later ``the model is parameterised by
$\boldsymbol\theta$.''
\end{notebox}

\begin{itemize}
  \item \kw{Decision stump} = depth-1 tree: enumerate all (feature, value) pairs, score each on the
  training set, return the best. The two parameters are \emph{which feature} and \emph{which value};
  the loss is implicit (error rate); the optimiser is brute force. It works here only because the
  parameter space is tiny ($n_{\text{features}}\times2$) — a deliberately \kw{weak} family that
  motivates the richer families and gradient-based optimisers of later weeks.
  \item \kw{Why ML beats hand-coded rules}: a hand rule must be rewritten from scratch for every new
  problem, so engineering effort scales with the number of problems; ML transfers the \emph{algorithm}
  and only swaps the data. \kw{Domain knowledge is not required} — the algorithm operates on the
  feature matrix and is blind to what the symbols mean, which is exactly why one algorithm serves
  fish-classification, fraud detection, and spam filtering alike.
  \item \kw{When to output a probability} rather than a hard label: when a downstream decision needs a
  \emph{confidence} to act on — fraud risk-pricing, medical risk, recommendation ranking. ``It is a
  fish'' and ``$p(\text{fish}\mid x)=0.83$'' answer different questions.
\end{itemize}

% =====================================================================
\section{Week 2 — Decision Trees}

A decision tree is the natural escalation of the Week~1 stump. A single stump puts one threshold on
one feature axis, which works only when the classes are separable by an \kw{axis-aligned} hyperplane.
Real data rarely cooperates — the canonical counter-example is a central disc of one class ringed by
the other, where every horizontal or vertical cut leaves half the ring on each side. The fix is to
apply the splitting procedure \emph{recursively}: split, then split each side again, and so on. The
decision boundary becomes a staircase of axis-aligned cuts that can approximate any curved boundary.

\begin{setupbox}[decision tree]
\kw{Internal node} = a test on one feature, routes data to a child. \kw{Leaf} = emits a prediction
(classification: majority class; regression: mean target). Training = \kw{recursive partitioning}:
at each node search candidate splits, pick the one most reducing impurity, recurse on children. The
model parameter $\boldsymbol\theta$ \emph{is} the whole tree — its topology, every node's
(feature, threshold), and every leaf's label — so unlike linear regression's fixed-shape weight
vector, the parameter object itself has variable shape.
\end{setupbox}

\begin{itemize}
  \item Training is \kw{greedy} — at each node it brute-forces the locally best single split, with no
  look-ahead to how that choice constrains descendants. Finding the globally optimal tree over all
  trees of bounded size is \kw{NP-hard}, so greedy is the standard tractable substitute: cheap, and
  good enough in practice.
  \item \kw{Linearly separable} = classes separable by a single hyperplane (a line in 2-D). Most real
  datasets are not linearly separable in their original feature space; trees sidestep this by
  stacking many axis-aligned splits rather than searching for one clever boundary.
\end{itemize}

\subsection{Splitting criteria}

The split criterion is the loss the tree minimises during training, and it must reward
\emph{partial} progress. Each node has class proportions $p_i$, summarised by an \kw{impurity}:
\[
H(p) = -\textstyle\sum_i p_i\log p_i \;\;\text{(entropy)}, \qquad
G(p) = 1-\textstyle\sum_i p_i^2 \;\;\text{(Gini impurity)}.
\]
Both are $0$ at a pure node and maximal at a uniform mix, and both vary \emph{smoothly} with the
proportions — which is the property that lets them detect a split that nudges a node toward purity
without yet flipping any prediction. Entropy is the average number of bits/nats to encode a sample
from the node's class distribution; Gini is the probability that two samples drawn with replacement
from the node have different classes. A split sends $n_l$ of the $n$ points left and $n_r$ right.

A split produces \emph{two} children, so a single comparable score is built by taking a
count-weighted average of the two child impurities — a branch carrying more data deserves more
weight, since it contributes more to the tree's future training set.

\begin{defbox}[split scoring — one template, two names]
\kw{Information Gain} (entropy), \emph{maximise}:
$\;I = H(p_{\text{par}}) - \big[\tfrac{n_l}{n}H(p_l)+\tfrac{n_r}{n}H(p_r)\big]$. \\
\kw{Gini split loss}, \emph{minimise}:
$\;L = \tfrac{n_l}{n}G(p_l)+\tfrac{n_r}{n}G(p_r)$. \\
Same structure = count-weighted average of child impurities; swap $H\leftrightarrow G$. Parent
impurity is constant at split time $\Rightarrow$ minimising the child sum $\equiv$ maximising the
gain. Information gain thus reads as ``how many bits of class uncertainty this split eliminates.''
\end{defbox}

For typical class proportions the entropy and Gini curves have the same shape — both vanish at the
pure endpoints, both peak at the even mixture — so the two criteria pick almost identical splits in
classification. Gini is the common default purely because it skips the logarithm.

\begin{trapbox}[IG and Gini interchangeable — only for classification]
Classification: both use the same $p_i$, behave near-identically. \kw{Regression}:
\begin{itemize}
  \item Entropy has a continuous counterpart $\Rightarrow$ IG becomes \kw{Variance Reduction}:
  $\text{Var}(\text{par})-\big[\tfrac{n_l}{n}\text{Var}(l)+\tfrac{n_r}{n}\text{Var}(r)\big]$
  (surrogate: Gaussian differential entropy $H=\tfrac12\log(2\pi e\sigma^2)$).
  \item Gini ($1-\sum p_i^2$) = discrete category-mismatch probability only; \kw{no} continuous form,
  no link to variance.
\end{itemize}
\end{trapbox}

For a \kw{regression tree}, the leaf prediction is the \kw{mean} of the training targets reaching
it, and the split criterion becomes \kw{variance reduction} — choose the split minimising the
count-weighted child variances $\tfrac{n_l}{n}\sigma_l^2+\tfrac{n_r}{n}\sigma_r^2$. Variance is the
regression analogue of Gini: zero on a pure node (all targets equal, so the mean is exactly right),
larger when the targets disagree. The \kw{median} is a more outlier-robust leaf prediction, since a
single huge target drags the mean far but the median only slightly.

\begin{pitfallbox}[0--1 loss is unusable \emph{inside} a tree]
0--1 loss is fine for a single \emph{stump} but breaks inside a tree because it is blind to ``less
wrong'': after the first split both children typically still hold a class mixture, and an
intermediate split can make them purer in \emph{distribution} without yet flipping any majority
prediction $\Rightarrow$ zero measured improvement $\Rightarrow$ the tree wastes splits and stalls.
The cure is an impurity measure (entropy/Gini) that varies smoothly with proportions and so rewards
distributional improvement at every level.
\end{pitfallbox}

\subsection{Input types}

The recursive-partitioning algorithm never changes across input types — only the list of
\emph{candidate splits} offered to the brute-force search does.

\begin{itemize}
  \item \kw{Real-valued}: sort the values along the axis; candidate thresholds sit \emph{between}
  adjacent sorted values (the midpoint convention keeps no training point exactly on the boundary);
  each ``$x<t$'' is one binary split.
  \item \kw{Quantised real} (pre-binned, e.g. age recorded as ``18--24''): the exact values have been
  discarded, so two points in the same bin are \kw{indistinguishable} to the model. A split
  \emph{inside} a bin therefore sends every point in it the same way and separates nothing (zero
  information gain) $\Rightarrow$ only test splits at \kw{bin boundaries}; the gain profile is a set
  of discrete spikes at the bin transitions.
  \item \kw{Categorical (unordered)}: there is no ``$<$'' to threshold, so use
  \kw{one-left-rest-right} — for $n$ categories test the $n$ splits ``$\{k\}\mid\text{rest}$'', each a
  boolean ``feature $==k$?''.
\end{itemize}

\begin{center}
\renewcommand{\arraystretch}{1.2}
\begin{tabularx}{\linewidth}{l l l X}
\toprule
\thead{Categorical strategy} & \thead{\# candidates} & \thead{Storage} & \thead{Scaling} \\
\midrule
All binary partitions & $O(2^{n-1})$ & subset of categories & exponential ($\sim5\times10^5$ at $n{=}20$) \\
One-left-rest-right & $O(n)$ & single category index & linear (20 at $n{=}20$) \\
\bottomrule
\end{tabularx}
\end{center}

After a categorical split both children grow further: left child = the $\{k\}$ rows (split on
\emph{other} features); right child = remaining-category mix (free to peel off another category).
The categorical feature is an \emph{input}, not the prediction target.

\begin{exbox}[what a categorical-feature split looks like]
Feature ``cheese'' $\in\{$Cheddar, Brie, Gouda, Stilton$\}$. A one-left-rest-right split on Brie:
\begin{center}\small\ttfamily
[\,cheese == Brie?\,]\quad yes $\to$ \{Brie rows\}\qquad no $\to$ \{Cheddar, Gouda, Stilton rows\}
\end{center}
The left child holds only the Brie rows and is grown on \emph{other} features (age, income, \dots);
the right child holds the remaining-category mixture and may split on ``cheese'' again to peel off
Cheddar, then Gouda, and so on. ``Cheese'' is an input feature being used to predict a separate
target — not the thing being predicted.
\end{exbox}

\subsection{Trade-offs}

\begin{pitfallbox}[interpretability $\ne$ accuracy]
Trees are \kw{interpretable} (readable decision path) — required in regulated domains (e.g. loan
decisions). But interpretability is a separate axis from accuracy: a readable tree is not necessarily
an accurate one.
\end{pitfallbox}

\begin{center}
\renewcommand{\arraystretch}{1.2}
\begin{tabularx}{\linewidth}{l l X}
\toprule
\thead{Train acc} & \thead{Test acc} & \thead{Diagnosis} \\
\midrule
high & low & \kw{overfitting} \\
low & low & \kw{underfitting} \\
high & high & good fit \\
low & high & data leakage / bug (suspicious) \\
\bottomrule
\end{tabularx}
\end{center}

% =====================================================================
\section{Week 3 — Bias, Variance, Cross-Validation}

Week~3 builds the language needed to say precisely \emph{why} a model is wrong. It starts with a
small worked case — the bias of the sample variance — purely to install the vocabulary ``bias of an
estimator = how far its expectation sits from the truth,'' then reuses that vocabulary in the
bias--variance decomposition.

\subsection{Naive sample variance is biased}

$\hat\sigma^2_{\text{naive}}=\tfrac1n\sum(X_i-\bar X)^2$ \kw{underestimates} $\sigma^2$. The reason
is that $\bar X$ is fitted to the very same data: $\sum(X_i-\bar X)^2$ is minimised at $\bar X$, so
deviations measured from the sample mean are systematically \emph{smaller} than deviations from the
true $\mu$. The estimator is reusing the data both to locate the centre and to measure spread about
it, and that double use is what costs a degree of freedom.

\begin{defbox}[sum-of-squares identity]
\[
\textstyle\sum_{i=1}^n(X_i-\bar X)^2 = \sum_{i=1}^n(X_i-\mu)^2 - n(\bar X-\mu)^2.
\]
\emph{Proof.} $X_i-\bar X=(X_i-\mu)-(\bar X-\mu)$; expand the square. Third term constant in $i$
$\to n(\bar X-\mu)^2$. Middle: $\sum(X_i-\mu)=n(\bar X-\mu)$ so it is $2n(\bar X-\mu)^2$. Combine
$-2n+n=-n$.\hfill$\blacksquare$
\end{defbox}

Taking expectations of the identity, $\mathbb E[\sum(X_i-\mu)^2]=n\sigma^2$ and
$\mathbb E[n(\bar X-\mu)^2]=n\,\text{Var}(\bar X)=\sigma^2$, so
$\mathbb E[\sum(X_i-\bar X)^2]=(n-1)\sigma^2$ — exactly $\sigma^2$ short of $n\sigma^2$. Hence the
\kw{unbiased estimator divides by $n-1$} (Bessel's correction): one degree of freedom was spent
estimating $\bar X$.

\subsection{Bias–variance decomposition}

The same ``how far is the expectation from the truth'' idea now applies to a trained model. A model
trained on a random dataset $\mathcal D$ is itself a \kw{random object}, written
$\hat f(\cdot;\mathcal D)$: re-sample $\mathcal D$ and you get a different fitted model. So we study
the \emph{expected} error, averaging over both the random training set and the test-point noise.

\begin{defbox}[MSE decomposition]
For target $y=f(x)+\varepsilon$, $\;\mathbb E[\varepsilon]=0$, $\;\text{Var}(\varepsilon)=\sigma^2$:
\[
\mathbb E_{\mathcal D,\varepsilon}\big[(y-\hat f(x))^2\big]
= \underbrace{(f(x)-\mathbb E_{\mathcal D}[\hat f(x)])^2}_{\text{Bias}^2}
+ \underbrace{\mathbb E_{\mathcal D}[(\hat f(x)-\mathbb E_{\mathcal D}[\hat f(x)])^2]}_{\text{Variance}}
+ \underbrace{\sigma^2}_{\text{irreducible}}.
\]
\end{defbox}

\begin{itemize}
  \item \examflag{}The \kw{proof is non-examinable}; the \kw{components, their meaning, and the
  decomposition idea are examinable}. Worth tracing the proof once: substitute $y=f+\varepsilon$,
  expand, drop the $\varepsilon$ cross term, then add-and-subtract $\mathbb E_{\mathcal D}[\hat f]$
  and expand again to peel off bias$^2$ and variance.
  \item \kw{Cross terms vanish} for two reasons, and articulating them is what shows you understand
  the proof: (i) $\varepsilon$ is independent of $\hat f$ — the noise on a test point has nothing to
  do with which training set produced the model — so $\mathbb E[(f-\hat f)\varepsilon]$ factorises and
  $\mathbb E[\varepsilon]=0$ kills it; (ii) $\mathbb E_{\mathcal D}[\mathbb E_{\mathcal D}[\hat f]-\hat f]=0$
  because the expectation of a quantity minus its own mean is zero.
  \item \kw{Dartboard} picture: train the model 100 times on 100 training sets and plot all 100
  predictions at one test point as darts. Bias = distance from the cluster's centre of mass to the
  bullseye (systematic offset); variance = how tightly the darts group around their own centre
  (consistency). The two are independent axes.
  \item \kw{Bias} = error from a model too rigid to represent the truth (a line fitting a circle —
  underfit); \kw{Variance} = error from a model flexible enough to track the sample's noise (overfit).
  A flexible family has low bias but high variance; a rigid family the reverse. The
  irreducible \kw{noise} $\sigma^2$ is the data's own ambiguity and no algorithm can touch it.
\end{itemize}

\begin{pitfallbox}[double descent]
The classical U-shaped trade-off is the examined picture, but it can \emph{break down} for very
over-parameterised models (deep learning) — test error descends again past the interpolation point.
\end{pitfallbox}

The trade-off is the lever behind almost every practical technique. Increasing model capacity
(deeper trees, higher-degree polynomials) lowers bias and raises variance; the total error
Bias$^2+\text{Variance}+\sigma^2$ is minimised at an \emph{intermediate} capacity, not at either
extreme. A fully grown tree sits at the low-bias, high-variance end — which is exactly why Week~4's
forests are principled rather than ad hoc: they average that variance away.

\subsection{Regularisation reduces variance}

\[
\text{Objective}=\text{Loss}(\hat f(X),Y)+\lambda\,\Omega(w), \qquad
\Omega=\textstyle\sum w_j^2\ \text{(L2/Ridge)}\ \text{or}\ \sum|w_j|\ \text{(L1/Lasso)}.
\]
\begin{itemize}
  \item The variance argument is concrete: a sharp wiggle threading a particular noisy point needs
  steep slopes, i.e. large weights — and those large weights are exactly what swings most from one
  training sample to the next, which \emph{is} the variance term. Penalising $\|w\|$ caps weight size,
  forbidding the model from contorting around any single point $\Rightarrow$ smoother predictions,
  less variance, a little bias added (the constrained optimum can no longer reach the unregularised
  one).
  \item $\lambda$ tunes the position on the bias--variance curve: $\lambda=0$ recovers the
  high-variance fit, larger $\lambda$ trades more bias for less variance. L1 additionally drives
  weights to \emph{exactly} zero (sparsity / feature selection), L2 only shrinks them.
\end{itemize}

\subsection{Data discipline and cross-validation}

Choosing where to sit on the bias--variance curve is an empirical decision, and it needs data the
model never trained on. That is the entire reason for splitting data three ways.

\begin{defbox}[three-way split]
\kw{Train} set — fits parameters. \kw{Validation} set — selects hyperparameters. \kw{Test} set —
used \emph{once}, final estimate only. Typical $\approx70$–$80\,/\,10$–$15\,/\,10$–$15$.
\end{defbox}

\begin{pitfallbox}[tuning on the test set]
Selecting hyperparameters on the test set leaks information $\Rightarrow$ optimistically biased
score. The test set must be touched once, at the end.
\end{pitfallbox}

A fixed train/validation split wastes data on small datasets — the validation points never train the
model. Cross-validation rotates the held-out slice through the data so every point is used for both.

\begin{defbox}[$k$-fold cross-validation]
Split data into $k$ equal folds. For each of $k$ rounds: hold out 1 fold (validate), train on the
other $k-1$, record score. Average the $k$ scores. Every point validates exactly once, and the
average smooths out the sampling noise of any single split.
\end{defbox}

\begin{itemize}
  \item Typical $k=5$ or $10$ (range $\sim4$–$20$); cost $=k\times$ training, since the model is
  refitted once per fold — with a grid search over hyperparameters this can dominate compute.
  \item Use: maximises usable data on small sets; averages out single-split sampling noise; it is the
  \kw{engine for hyperparameter tuning}. The full pipeline: pick the best hyperparameter by CV
  $\to$ retrain one final model on all the train+validation data $\to$ touch the test set once.
  \item \kw{Stratified} $k$-fold = preserves class proportions per fold (imbalanced classes).
  \kw{LOOCV} = $k{=}N$: thorough, $N$ retrains, expensive.
\end{itemize}

% =====================================================================
\section{Week 4a — Random Forests}

Week~3 diagnosed the single tree's weakness precisely — low bias, high variance — and the natural
cure is to \emph{average many trees} so their random fluctuations cancel. That only works if the
trees genuinely \kw{disagree}: a hundred identical trees average to one tree. So the whole design of
a random forest is machinery for manufacturing \emph{diverse} trees from one fixed dataset.

\begin{defbox}[bootstrap sample]
From a training set of size $n$: draw one point uniformly \kw{with replacement}, repeat exactly $n$
times $\Rightarrow$ a same-size-$n$ sample with duplicates and omissions. Repeat $S$ times for $S$
datasets (one per tree). Bootstrapping is the cheap way to ``fake'' the many training sets we cannot
collect — each draw is a slightly different mix of the same data.
\end{defbox}

\begin{itemize}
  \item Two loops: the \kw{inner} loop samples $n$ times to build one dataset; the \kw{outer} loop
  runs $S$ times, one dataset per tree.
  \item A specific point survives a single draw with probability $1-\tfrac1n$, so survives all $n$
  draws with probability $(1-\tfrac1n)^n\to 1/e\approx0.37$. Hence each bootstrap sample omits
  $\approx37\%$ of the original points — that omitted third is the tree's \kw{out-of-bag (OOB)} set,
  a free held-out set for \emph{that} tree.
\end{itemize}

Bagging alone is not quite enough for trees: the same few dominant features keep winning the root
split, so bagged trees still look alike. The random subspace method forces the issue.

\begin{defbox}[random subspace method]
At \kw{every split} (re-sampled per split, not once per tree), restrict the candidate features to a
fresh random subset. Subset size: $\sqrt n$ (classification), $n/3$ (regression), of $n$ total
features. By denying trees their favourite features it pushes them to base early decisions on
different features, decorrelating them further.
\end{defbox}

\begin{itemize}
  \item \kw{Bagging} = \kw{B}ootstrap \kw{agg}regat\kw{ing} — the resampling is on the data, but what
  gets aggregated is the prediction of each \emph{fully trained} estimator. \kw{Final prediction}:
  push a new point through \emph{all} $S$ trees; average (regression) / majority vote (classification).
  \item \kw{OOB evaluation} gives a free validation set: to score training point $i$, aggregate only
  the trees whose bootstrap sample \emph{omitted} $i$. No tree is ever judged on a point it trained
  on, so the OOB error is an honest generalisation estimate at zero extra cost — and it is a distinct
  operation from final prediction, which uses every tree.
  \item Bootstrap + random subspace $\Rightarrow$ \kw{diverse} trees $\Rightarrow$ near-independent
  errors. Averaging $S$ independent estimators each of variance $V$ gives variance $V/S$ (the
  variance of a mean of $S$ i.i.d.\ quantities) $\Rightarrow$ \kw{bias preserved} (each tree is still
  flexible and trained on data from the true distribution, so the average stays low-bias),
  \kw{variance destroyed}. This is the whole point: the forest improves the variance term without
  paying in bias.
  \item \kw{Boosting} = the other ensemble family: trees built \kw{sequentially}, each correcting the
  errors of those before it (vs bagging's \kw{parallel}, independent trees).
  \item \kw{Limitations}: a forest sacrifices the single tree's \kw{interpretability} (you can read
  one tree, not 100 averaged ones); like any tree it \kw{cannot extrapolate} beyond the range of the
  training data; and it costs more memory and prediction time (all $S$ trees stored and evaluated).
\end{itemize}

\begin{notebox}[why ensembling works, and when to use a forest]
\kw{Why it works}: on noisy or limited data many different trees fit the training set about equally
well — committing to a single one throws information away. Averaging a \emph{diverse} set keeps the
shared signal and lets the independent errors cancel. \kw{Bagging alone is not enough}: bootstrap
samples still share the dominant features, so the trees stay correlated — the random subspace method
is what forces genuine diversity. \kw{When}: choose a forest for raw accuracy (pre-deep-learning it
was the best general-purpose classifier, and it is fast to train); choose a single tree when you
need interpretability — you cannot have both at once.
\end{notebox}

% =====================================================================
\section{Week 4b — Loss \& Performance Measures}

\subsection{Risk and the Bayes classifier}

This half of Week~4 switches register from algorithms to evaluation: what does it actually mean for a
classifier to ``perform well''? The answer needs a clean separation between the error we can measure
and the error we genuinely care about.

\begin{defbox}[empirical vs expected risk]
\kw{Empirical risk} $R_{\text{emp}}=\tfrac1n\sum_i\ell(\hat f(x_i),y_i)$ — the average loss on a
finite \emph{sample}; measurable. \\
\kw{Expected (true) risk} $R(f)=\mathbb E_P[\ell(f(x),y)]$ — the average loss over the true
distribution $P$; what we want, but \emph{not} measurable since $P$ is unknown. A held-out test set's
empirical risk is an unbiased estimate of it, provided the test data really is drawn from $P$ and was
never seen during training.
\end{defbox}

\begin{defbox}[Bayes classifier]
Predicts, for each $x$, the class with the highest \kw{true} posterior $P(C_k\mid x)$ — for binary
classes, predict $C_1$ when $P(C_1\mid x)\ge0.5$. Predicting the more probable class is provably the
lowest-error rule, so it is the optimum. But it assumes $P$ is known exactly, which it never is
$\Rightarrow$ the Bayes classifier is a \kw{theoretical benchmark} you can reason about but never
build.
\end{defbox}

\begin{trapbox}[knowing the distribution removes uncertainty — FALSE]
Knowing $P$ gives the exact \emph{odds}, not the \emph{outcome}: $x\mapsto P(y\mid x)$ is
\kw{stochastic}. The optimal classifier's unavoidable error = \kw{Bayes error rate} = irreducible
\kw{noise floor} = \kw{aleatoric uncertainty} (hidden variables + intrinsic randomness).
\[
\text{model error} = \underbrace{\text{Bayes error}}_{\text{noise floor}}
+ \underbrace{\text{excess risk}}_{\text{estimator's fault}}.
\]
\end{trapbox}

\subsection{Conditional probability from a joint density}

To reach the posterior the Bayes classifier needs, two standard rules combine: the definition of
conditional probability $P(C_1\mid x)=p(x,C_1)/p(x)$, and the law of total probability for the
marginal $p(x)=p(x,C_1)+p(x,C_2)$.

\kw{Notation}: $P(\cdot)$ = probability (a pure number in $[0,1]$); $p(\cdot)$ = probability
\kw{density} (the height of a curve, which can exceed $1$ — only $p\cdot\mathrm dx$ is a probability).
\[
P(y{=}C_1\mid x)=\frac{p(x,C_1)}{p(x)},\qquad p(x)=p(x,C_1)+p(x,C_2)
\;\Rightarrow\;
P(y{=}C_1\mid x)=\frac{p(x,C_1)}{p(x,C_1)+p(x,C_2)}.
\]
Bayes' rule for continuous variables. The window factor $\mathrm dx$ ($P=p\cdot\mathrm dx$) is
identical in every term and \kw{cancels} $\Rightarrow$ ratio of curve heights; result on the left is
a genuine probability.

\subsection{Bayes-optimal threshold}

\begin{setupbox}[1-D two-class threshold]
Two joint-density curves $p(x,C_1),p(x,C_2)$ over an axis. \kw{Crossover $x_0$} = curves equal
$\Rightarrow P(C_1\mid x_0)=\tfrac12$ = \kw{Bayes-optimal threshold}. It slices the axis into
\kw{decision regions} $R_1$ (predict $C_1$, taller curve), $R_2$ (predict $C_2$).
\end{setupbox}

\begin{itemize}
  \item Classification is a property of \kw{regions}, not point values. \kw{Errors are area
  integrals}: Bayes error $=\int_{R_1}p(x,C_2)\,\mathrm dx+\int_{R_2}p(x,C_1)\,\mathrm dx$ (the
  unavoidable curve overlap).
  \item Density height at a point is not a probability; the height-ratio is the conditional
  probability of a \kw{micro-region} of width $\mathrm dx$ — string these together to build the
  macro-regions.
  \item \kw{$\hat x$} = the \emph{estimated} boundary of a real trained model, $\hat x\ne x_0$. In the
  band between $x_0$ and $\hat x$ the model predicts the shorter-curve (wrong) class $\Rightarrow$ the
  \kw{red region} = \kw{excess risk}, the avoidable error from the estimator. Green/blue overlap =
  the universe's fault; red = the estimator's fault.
\end{itemize}

\begin{exbox}[Bayes-optimal threshold — worked numbers]
At a point $\hat x$ suppose the class-density heights are $p(\hat x,C_1)=1$ and $p(\hat x,C_2)=4$. The
conditional probabilities are
\[
P(C_1\mid\hat x)=\frac{1}{1+4}=0.20,\qquad P(C_2\mid\hat x)=\frac{4}{1+4}=0.80
\]
— they sum to $1$, and the Bayes rule predicts the larger, $C_2$. At the crossover $x_0$ the two
heights are equal ($p$ each), so $P(C_1\mid x_0)=\tfrac{p}{2p}=0.5$ — exactly the decision boundary.
A trained model whose boundary $\hat x$ sits away from $x_0$ predicts the $20\%$ class throughout the
band between them: that band is the avoidable excess-risk error.
\end{exbox}

\subsection{0--1 loss, margin, surrogates}

0--1 loss is the right \emph{metric} — its empirical version is exactly the classification error — but
the wrong \emph{loss function} to optimise. The two roles must be kept apart: the metric is what you
report and the business cares about; the training loss is a tractable surrogate the optimiser can
actually descend.

Binary labels $y\in\{-1,+1\}$, raw score $f(x)$. \kw{Margin} $z=y\,f(x)$ is positive when $y$ and
$f(x)$ share a sign ($z>0$ correct, $z\le0$ wrong). \kw{0--1 loss} as a function of $z$ is a step:
$1$ for $z\le0$, $0$ for $z>0$.

\begin{notebox}[why 0--1 loss is unusable for gradient methods]
\kw{Zero gradient} almost everywhere, undefined at $z{=}0$ $\Rightarrow$ no descent direction.
Summed over a training set: \kw{non-convex}, combinatorially many local minima, exact minimisation is
\kw{NP-hard}. Empirically a multi-dimensional staircase of flat plateaus + cliffs.
\end{notebox}

\kw{Convex surrogates} replace the step with a tractable upper bound: convex (so a global minimum is
reachable), differentiable (so gradient methods have a direction), and lying \emph{above} the step
(so minimising the surrogate also approximately minimises 0--1 loss). All three below have a usable
slope on the wrong side $z<0$:
\begin{itemize}
  \item \kw{Hinge} $\max(0,1-z)$ (SVM) — zero once the margin is comfortable ($z\ge1$), linear
  penalty below; it still penalises correct-but-unconfident points ($0<z<1$) but stops pushing once
  the margin is large, a useful sparsity-inducing property.
  \item \kw{Logistic} $\log(1+e^{-z})$ — smooth, strictly positive, decays towards 0 but never
  reaches it, so it always pushes for a little more confidence.
  \item \kw{Squared} $(1-z)^2$ — convex but \kw{increases again for $z\gg1$}, so it over-penalises
  confident-correct predictions $\Rightarrow$ suboptimal for classification (it is the natural choice
  for regression instead).
\end{itemize}

\subsection{Confusion matrix and metrics}

\begin{defbox}[confusion matrix — this unit's convention]
\kw{Rows = predicted, columns = actual} (some textbooks transpose — check the axes). \\
Read each cell backwards: 2nd word = predicted class; 1st word = whether correct.
\kw{TP} predicted +, correct. \kw{TN} predicted $-$, correct. \kw{FN} predicted $-$, wrong (the
miss — actually +). \kw{FP} predicted +, wrong (false alarm — actually $-$).
\end{defbox}

\[
\text{Acc}=\frac{TP+TN}{TP+TN+FP+FN},\quad
\text{Prec}=\frac{TP}{TP+FP},\quad
\text{Rec}=\frac{TP}{TP+FN},\quad
F_1=\frac{2\,\text{Prec}\cdot\text{Rec}}{\text{Prec}+\text{Rec}}.
\]

Read the metrics by the question each answers. \kw{Precision} = ``of the things I called positive,
how many actually were?'' — high precision means few false alarms (a spam filter optimises this).
\kw{Recall} = ``of the actual positives, how many did I catch?'' — high recall means few misses
(cancer screening optimises this). The two trade off: always-predict-positive maxes recall and tanks
precision. \kw{$F_1$} is their \emph{harmonic} mean, which sits closer to the smaller of the two, so
a high $F_1$ forces \emph{both} to be good — you cannot game it by acing one and tanking the other.

\begin{trapbox}[balanced accuracy = \emph{mean} of recalls]
Not the sum (could reach 2.0). Binary: $\tfrac12(\text{Rec}_++\text{Rec}_-)$. Multi-class: sum
per-class recalls, divide by \#classes.
\end{trapbox}

\begin{pitfallbox}[accuracy misleads on imbalanced data]
99.9\%-negative fraud data: ``always predict negative'' scores 99.9\% accuracy yet is useless. Use
\kw{F1} or \kw{balanced accuracy} instead.
\end{pitfallbox}

\begin{exambox}[Q5-style — compute metrics from $TP{=}40,\ FN{=}10,\ FP{=}20,\ TN{=}30$]
Prec $=40/60=0.667$; Rec $=40/50=0.80$; Acc $=70/100=0.70$;
$F_1=2(0.667)(0.8)/(1.467)=0.727$; BalAcc $=\tfrac12(0.80+0.60)=0.70$
(Rec$_-=TN/(TN+FP)=30/50=0.60$).
\end{exambox}

\begin{center}\large\bfseries\color{defBd} PART II — Optimisation (Weeks 7–9)\end{center}

% =====================================================================
\section{Week 7 — Optimisation framing}

Model fitting \emph{is} optimisation — pick a family, pick a loss, minimise. Week~7 supplies the
language and the toolkit for that minimisation step; Weeks~8 and~9 instantiate it for linear
regression with gradient descent and Newton's method. The framing example throughout is ``find the
biggest apple on the tree'': with one attribute the search space is a smooth single bowl a computer
solves trivially; add attributes and the surface becomes high-dimensional and bumpy. The asymmetry
that drives everything is that the optimiser, unlike a human eyeing a 3-D plot, sees only the
function value at points it has explicitly evaluated — it must decide where to look next from a
local view.

\begin{defbox}[gradient and Hessian]
\kw{Gradient} $\nabla f$ = column vector of first partial derivatives. It points in the direction of
fastest \emph{increase}, and its magnitude is the rate of that increase — so it carries both ``which
way'' and ``how steep.'' \kw{Hessian} $\nabla^2 f$ = matrix of second partials; symmetric (Schwarz's
theorem, for $C^2$ functions). It encodes \kw{curvature} — how fast the slope itself is changing —
and is what Week~9's Newton's method consumes.
\end{defbox}

\subsection{The generic iterative scheme}

This is the single most examinable abstraction in Week~7: every gradient-based optimiser in the unit
is just a specific choice of the three design knobs below.

\begin{defbox}[iterative optimiser]
Initialise $x_0$; repeat $x_{t+1}=x_t+\alpha_t\,p_t$ until termination. Three design choices:
\kw{direction} $p_t$ (gradient descent: $p_t=-\nabla f(x_t)$ — the negative gradient, because the
gradient points uphill and we want downhill), \kw{step size} $\alpha_t$ (the learning rate; too
large overshoots and oscillates, too small crawls), and the \kw{termination condition}.
\end{defbox}

\kw{Termination conditions} (any one):
\begin{itemize}
  \item Iteration cap $t\ge T_{\max}$.
  \item Gradient near zero $\|\nabla f(x_t)\|\le\epsilon_g$.
  \item Progress near zero $\|x_{t+1}-x_t\|\le\epsilon_x$.
\end{itemize}

\begin{pitfallbox}[$t\ge T_{\max}$ is the \emph{stop} rule]
The notes state the \emph{termination} (stop) rule: from $t{=}0$, stop when $t\ge T_{\max}$. The
familiar $t<T_{\max}$ is the \emph{continuation} (while-loop) condition — the logical negation.
\end{pitfallbox}

\begin{pitfallbox}[stopped moving $\ne$ found the minimum]
GD converges to a \kw{stationary point} ($\nabla f=0$) — could be a local min, saddle, or plateau.
Only \kw{convexity} licenses ``stationary $\Rightarrow$ global minimum''.
\end{pitfallbox}

\subsection{Convexity}

Convexity is the property that licenses an optimiser's local view. A function is \kw{convex} if every
chord joining two points on its graph lies on or above the graph; the consequence is that
\emph{every local minimum is the global minimum} — the bowl has one bottom. So any descent method
that reaches a stationary point is provably correct. Strict convexity (chord strictly above) further
makes the optimum \emph{unique}. Most non-trivial ML losses (deep networks, mixtures, hidden-variable
models) are non-convex, with many basins; from inside one there is no way to tell a better one
exists. The working strategy is then ``perform convex within non-convex'': treat each local region
as if convex and accept the local optimum it yields.

\begin{center}
\renewcommand{\arraystretch}{1.25}
\begin{tabularx}{\linewidth}{l X X}
\toprule
\thead{Aspect} & \thead{Convex} & \thead{Non-convex} \\
\midrule
Local minimum & is the global minimum & may be a worse-than-global trap \\
Initial guess matters? & No — same answer each run & Yes — different start, different answer \\
Theory & clean guarantees & no general theory; heuristics \\
Strategy & run any descent method to convergence & split surface, treat regions as convex, accept local optimum \\
\bottomrule
\end{tabularx}
\end{center}

\begin{pitfallbox}[convexity is about the loss surface, not the data]
Linear and logistic regression have convex loss \kw{regardless} of how the data is distributed.
Convexity is a property of $L(w)$.
\end{pitfallbox}

\subsection{Four families of optimisers}

The four families form a spectrum from guaranteed-but-unscalable to scalable-but-local. The list
itself is examinable.

\begin{enumerate}
  \item \kw{Brute force} / exhaustive — try every candidate, return the best. Guaranteed optimal but
  blows up combinatorially ($N$-queens reaches $64^{64}\!\approx\!10^{115}$); only for tiny discrete
  spaces.
  \item \kw{Random / grid search} — sample the space rather than exhaust it; cheaper, no optimality
  guarantee. \kw{Random beats grid} when only some inputs matter: a grid wastes evaluations re-testing
  the same value of the important axis, whereas random sampling gives a fresh value of every axis per
  point, probing the axis that matters far more densely (Bergstra \& Bengio).
  \item \kw{Metaheuristics} — population or stochastic search that tolerates uphill moves to escape
  local minima.
  \item \kw{Gradient-based} — use $\nabla f$ (and optionally $\nabla^2 f$) to pick a guaranteed-descent
  direction; fast, principled, scales to high dimensions, but needs a differentiable $f$ and gets
  stuck in local minima on non-convex losses. The workhorse of the rest of the unit.
\end{enumerate}

\begin{notebox}[the core challenge — an optimiser only sees locally]
A gradient-based optimiser never sees the whole loss surface; at each step it sees only a small
neighbourhood and must decide where to look next. That local-only view is \emph{why} convexity
matters so much (it guarantees the local ``downhill'' direction is globally right) and \emph{why}
non-convex problems need heuristics — from inside one basin there is no way to know a better basin
exists elsewhere.
\end{notebox}

\subsection{Metaheuristics}

These are controlled-randomness algorithms borrowed from physical or biological processes; you should
be able to describe each in a sentence or two, not derive them.

\begin{itemize}
  \item \kw{Particle swarm} — a population of candidates scattered through the space; each moves by a
  blend of local descent, its own best-so-far, and the swarm's shared global best. Balances
  \kw{exploration} and \kw{exploitation}: a particle stuck in a shallow basin is pulled out by
  another particle's discovery, so collective shared information outperforms a single greedy walk.
  \item \kw{Simulated annealing} — mimics heating then slowly cooling a metal. It maintains a
  candidate and accepts uphill steps with a temperature-controlled probability; early (hot) iterations
  jump widely and tolerate uphill moves to escape local minima, then it \kw{cools} so late iterations
  only accept downhill moves and freeze into a basin.
  \item \kw{MCMC} — walks a memoryless Markov chain (next state depends only on the current one)
  whose long-run visit frequencies match an unknown target distribution; the histogram of visits
  $\approx$ the distribution. It is a \kw{sampler}, not an optimiser — it maps the whole landscape and
  is used for Bayesian inference. Same mechanism as simulated annealing but \kw{no cooling}. Slow:
  millions of evaluations, so run multiple chains in parallel.
\end{itemize}

\begin{trapbox}[``Newton's method applies to any optimisation problem'' — FALSE]
Newton needs $f$ twice-differentiable and the Hessian invertible. See §10 for the applicability test.
\end{trapbox}

% =====================================================================
\section{Week 8 — Gradient Descent}

Week~8 plugs concrete content into Week~7's skeleton on the workhorse problem of \kw{linear
least-squares regression}, and solves it two ways — a closed form and an iterative descent.

\begin{defbox}[convex set and convex function]
\kw{Convex set} $C$: for all $x,y\in C$, $\beta\in[0,1]$, $\;(1-\beta)x+\beta y\in C$ — every line
segment joining two points of $C$ stays inside (no dents, holes, or separated pieces).
\kw{Convex function}: its domain is a convex set and the chord lies on/above the graph; equivalently
its Hessian is positive semi-definite. The headline property is the one from Week~7 — any local
minimum is the global minimum, and uniquely so under strict convexity.
\end{defbox}

\subsection{Batch GD}

Batch GD uses the whole dataset per step. The total loss is a \emph{sum} of per-point losses, and
the derivative of a sum is the sum of derivatives $\Rightarrow$ the batch gradient is the average of
the per-point gradients — nothing is approximated, this is the exact gradient. The point is that you
never run a $\sum_i$ loop: the whole sum collapses into one matrix expression evaluated in a single
BLAS call. (Squared error is preferred over absolute/$L_1$ error because it is \kw{differentiable
everywhere} — a clean gradient and a closed form — whereas $|\cdot|$ has a non-differentiable kink at
$0$.)
\begin{setupbox}[regularised least squares]
\[
L(w)=\|Xw-y\|^2+\lambda\|w\|^2, \qquad
\nabla L(w)=2X^\top(Xw-y)+2\lambda w.
\]
The product $X^\top(\cdots)$ \emph{is} the engine summing every per-point gradient at once.
\end{setupbox}

\kw{Normal equation} — the first solution method. Since the loss is convex, its unique stationary
point is the global minimum, so set the (unregularised) gradient to zero and solve:
$\nabla L=2X^\top(Xw-y)=0\Rightarrow X^\top X\,w^\ast=X^\top y$, giving
$w^\ast=(X^\top X)^{-1}X^\top y$ \kw{if $X^\top X$ is invertible} ($X$ full column rank). It is an
exact answer in one shot, a single line of NumPy — but it costs an $M\times M$ inverse, which is
$O(M^3)$ and infeasible for large $M$.
\begin{itemize}
  \item Why is the loss convex at all? The Hessian is $2X^\top X$, and
  $z^\top X^\top X\,z=\|Xz\|^2\ge0$ for every $z$ — a squared norm is never negative — so $X^\top X$
  is PSD \emph{whatever} the data. Invertible $\Rightarrow$ Hessian PD $\Rightarrow$ strictly convex
  $\Rightarrow$ \kw{unique} $w^\ast$.
  \item Singular $X^\top X$ (collinear features, or fewer points than features) $\Rightarrow$ Hessian
  only PSD $\Rightarrow$ a flat-bottomed valley, so \kw{infinitely many} $w^\ast$ all achieve the
  minimum.
\end{itemize}

\begin{trapbox}[cannot split $(X^\top X)^{-1}=X^{-1}(X^\top)^{-1}$]
$(AB)^{-1}=B^{-1}A^{-1}$ needs $A,B$ \kw{square}. Design matrix $X$ is $N\times M$, $N\gg M$,
rectangular $\Rightarrow$ $X^{-1}$ does not exist. Only the square product $X^\top X$ ($M\times M$)
is invertible.
\end{trapbox}

\begin{pitfallbox}[do not literally invert the Gram matrix]
In practice solve the linear system $X^\top X\,w=X^\top y$ (Cholesky / QR), not by forming
$(X^\top X)^{-1}$ — faster and numerically safer.
\end{pitfallbox}

\subsection{Near-collinearity / ill-conditioning}

Highly correlated columns carry near-identical information. Determinant = volume spanned by columns;
near-parallel column vectors collapse that volume $\Rightarrow\det(X^\top X)\to0$, matrix
\kw{ill-conditioned}. Computing $(X^\top X)^{-1}$ divides by the near-zero determinant $\Rightarrow$
inverse entries \kw{blow up} $\Rightarrow$ tiny floating-point rounding errors are \kw{amplified}
until fitted weights are noise, not signal. Cure: ridge $+\lambda I$ lifts the determinant off zero.

\subsection{SGD and mini-batch}

Batch GD's exact gradient costs $O(NM)$ per step — slow when $N$ is huge. The variants trade
gradient accuracy for cheaper steps by averaging over fewer points.

\begin{defbox}[SGD update]
SGD abandons the full matrix $X$ entirely: draw one random index $i_s$ and step on that single
point's noisy gradient. For least squares $\nabla L_i=2(w^\top x_i-y_i)x_i$, so
\[
w_{t+1}=w_t-2\alpha\,(w_t^\top x_{i_s}-y_{i_s})\,x_{i_s}.
\]
\end{defbox}

\begin{defbox}[mini-batch GD]
Batch of $B$ points (powers of 2 for GPUs): $\;w_{t+1}=w_t-\dfrac{\alpha}{B}\sum_{i\in\text{batch}}\nabla L_i$.
Lower-variance gradient estimate than SGD. ``SGD'' in modern usage = mini-batch.
\end{defbox}

The three variants sit on a spectrum from expensive-and-stable to cheap-and-noisy:
\begin{itemize}
  \item Batch GD: slow, exact-gradient steps that head straight to the minimum with no sampling noise.
  \item SGD: fast and cheap, but the single-point gradient is noisy, so the trajectory is a
  ``drunken'' zig-zag; on average it still descends, and faster in wall-clock terms. \kw{Bonus}: the
  noise can \emph{escape} narrow local minima on non-convex losses — an advantage for deep learning.
  \item Mini-batch: the practical compromise — a small matrix of 32/64/128 rows gives a
  lower-variance gradient than SGD while staying far cheaper than batch. ``SGD'' in modern usage means
  mini-batch.
\end{itemize}

\begin{exambox}[Q3-style T/F — ``convex + right step size $\Rightarrow$ BGD reaches the global optimum, SGD not''. \trapflag]
\kw{TRUE as stated.} Convex = single global min; BGD uses the exact full-data gradient $\Rightarrow$
deterministic descent to the bottom — ironclad. SGD's single-point gradient is noisy: fixed step
$\Rightarrow$ bounces forever; converges only ``in expectation'' with a \kw{decaying} step schedule —
no unconditional guarantee. ``Guaranteed to eventually converge'' $\Rightarrow$ mark TRUE.
\end{exambox}

\subsection{Anisotropy — why GD zig-zags}

How well steepest descent performs depends on the \emph{shape} of the loss surface, not just its
convexity.

\begin{itemize}
  \item \kw{Isotropic} loss (circular contours) $\Rightarrow$ $-\nabla f$ — which is always
  perpendicular to the local contour — happens to point straight at the minimum, so descent is
  direct.
  \item \kw{Anisotropic} loss (elliptical contours, large \kw{condition number} = ratio of largest
  to smallest Hessian eigenvalue) $\Rightarrow$ perpendicular-to-contour is \emph{not} towards the
  minimum: the iterate overshoots the long axis, is corrected, overshoots again — the slow
  \kw{zig-zag}. In the extreme it bounces across a thin valley with almost no forward progress.
  Newton's method fixes this by rescaling each direction with the local curvature.
  \item \kw{Trap}: anisotropy is \emph{not} non-convexity — a badly-conditioned quadratic is still
  perfectly convex; it has one global minimum, it just makes GD zig-zag on the way there.
\end{itemize}

% =====================================================================
\section{Week 9 — Newton's Method}

Newton's method closes the optimisation arc with a \emph{second-order} method: where gradient descent
uses only $\nabla f$ and a fixed scalar step, Newton's also reads the \kw{curvature} off the Hessian
and uses it to set a per-direction step automatically. This is exactly what cures GD's
anisotropy zig-zag.

\begin{defbox}[curvature]
Curvature = rate of change of the slope = $f''$ (1-D) / Hessian (multivariate). Geometrically it is
$1/r$ where $r$ is the radius of the inscribed circle that best fits the curve: a tight turn means a
small circle and \emph{high} curvature, a flat stretch means a huge circle and near-zero curvature.
The optimiser cares because near a sharp turn it should take a small step (or overshoot) and on a
flat stretch a large one — and GD's fixed $\alpha$ cannot tell the two apart.
\end{defbox}

\subsection{Derivation and update}

The update is derived by approximating $f$ near the current iterate with a quadratic, then jumping
straight to that quadratic's minimum. The 2nd-order Taylor model at $w_t$ is
$\;f(w)\approx f(w_t)+\nabla f^\top(w-w_t)+\tfrac12(w-w_t)^\top H(w-w_t)$. Differentiate w.r.t.\ $w$
to get $\nabla f(w_t)+H(w_t)(w-w_t)$, set this to zero, and solve $\Rightarrow$
\begin{defbox}[Newton update]
\[
w_{t+1}=w_t-[H(w_t)]^{-1}\nabla f(w_t).
\]
Fits a local quadratic, jumps to its minimum. There is no separate step size: the matrix $H^{-1}$
\emph{is} the step size, generalised from a scalar — it rescales every direction by its own
curvature, which is why on an anisotropic bowl it ``rounds out'' the contours and points straight at
the minimum. Requires $H$ \kw{invertible}.
\end{defbox}

\subsection{Convergence behaviour}

Whether Newton lands in one step depends entirely on whether the loss is exactly quadratic.

\begin{itemize}
  \item \kw{Quadratic loss}: $H=2X^\top X$ is \emph{constant} (two differentiations of
  $(w^\top x_i-y_i)^2$ kill the $w$). Substituting into the update, the $\tfrac12$ from $H^{-1}$ and
  the $2$ from the gradient cancel, and $(X^\top X)^{-1}(X^\top X)=I$ cancels the leading $w_t$:
  $\;w_{t+1}=(X^\top X)^{-1}X^\top y$. Both the starting point and the iteration index vanish — Newton
  hits the exact normal-equation solution in \kw{one step}, from anywhere. The reason is that the
  2nd-order Taylor model of a quadratic has \emph{zero remainder}: the local model \emph{is} the
  global loss.
  \item \kw{Higher-order loss} (e.g. $\sum(\cdots)^4$): two differentiations drop the power from 4 to
  2 but not to 0, so $H$ still contains $w$ — it is \kw{$w$-dependent}, curvature varies with
  position $\Rightarrow$ the local parabola matches only locally $\Rightarrow$ Newton must
  \kw{iterate}. Even so the rate is \kw{quadratic} near the optimum (the number of correct digits
  roughly doubles per step) — far faster than GD's merely \kw{linear} rate.
\end{itemize}

\begin{notebox}[Taylor expansion to the loss's own degree]
A degree-$N$ polynomial loss expanded to Taylor order $N$ has zero remainder — the model equals the
true loss everywhere $\Rightarrow$ one-shot convergence recovered. We stop at order 2 because order
$\ge3$ needs derivative \kw{tensors} ($M^3,M^4,\dots$ — storage explodes) and minimising a
cubic/quartic local model is no longer a fast linear solve.
\end{notebox}

\subsection{Applicability}

\begin{defbox}[two-line applicability test]
Newton applies iff: (1) $f$ is \kw{twice differentiable} ($C^2$); (2) the Hessian is
\kw{invertible} at every iterate.
\end{defbox}

\begin{notebox}[loss must leave $w$ after first differentiation — wording correction]
If $g(w)$ is at most first-order (linear), $\nabla g$ is constant and $H=0$ $\Rightarrow$ $H^{-1}$
does not exist, Newton cannot run. Quadratic $\Rightarrow$ $H$ constant non-zero $\Rightarrow$ one
step. Higher-order $\Rightarrow$ $H$ depends on $w$ $\Rightarrow$ iterate. \\
The notes' phrase ``\emph{must not collapse to a constant under second differentiation}'' is
misleading — a non-zero constant Hessian is the \emph{ideal}. Accurate: ``must not collapse to a
constant under \kw{first} differentiation'' ($\equiv$ not to \kw{zero} under second).
\end{notebox}

\subsection{Three failure modes}

Each failure mode is one of the applicability conditions being violated.

\begin{enumerate}
  \item \kw{Singular Hessian} (zero curvature in some direction = zero eigenvalue) $\Rightarrow$ the
  step is \kw{undefined} — note: undefined, not infinite. A flat region has no parabola, so the
  question ``where is the local model's minimum?'' has no answer at all.
  \item \kw{Not twice-differentiable} (kinks: hinge, $\ell_1$, ReLU) $\Rightarrow$ no Hessian exists
  at the kink — use sub-gradient methods or smooth the kink locally.
  \item \kw{Initial point far from optimum} — the 2nd-order Taylor model is only \emph{locally}
  accurate, so from far away a full Newton step can shoot off into a worse region; and on a
  non-convex surface Newton heads to whatever \kw{local} minimum is nearest, with no way to know a
  better basin exists. Newton has \emph{local} convergence guarantees only.
\end{enumerate}

\subsection{Damping (fix for singular / unreliable Hessian)}

There are two distinct kinds of damping for two distinct problems. Shrinking the step with a scalar
tames \emph{magnitude} (failure mode~3) but does nothing for a \emph{singular} Hessian (failure
mode~1) — there $H^{-1}$ still does not exist, so the damping must go \emph{inside} the inverse.

\kw{Levenberg–Marquardt}: use $(H+\lambda I)^{-1}$. Adding $\lambda I$ shifts \kw{every} eigenvalue by
$\lambda$: $e_i\to e_i+\lambda$. A zero eigenvalue becomes $\lambda\ne0$, restoring invertibility; for
the damped matrix to be \emph{positive definite} as well you need every shifted eigenvalue $>0$,
i.e.\ $\lambda>-e_{\min}$.

\kw{Choosing $\lambda$}:
\begin{itemize}
  \item \kw{Levenberg–Marquardt} (dynamic): step succeeds $\to$ shrink $\lambda$ (trust Hessian, $\to$
  pure Newton); step fails $\to$ grow $\lambda$ ($\to$ small-step gradient descent).
  \item \kw{Trust-region}: fix a radius the quadratic model is trusted within; $\lambda$ emerges as
  the Lagrange multiplier enforcing it.
  \item \kw{Static} (ridge): fixed hyperparameter, chosen by cross-validation.
\end{itemize}

\kw{Damped Newton}: $w_{t+1}=w_t-\alpha_t H^{-1}\nabla f$, $\alpha_t\in(0,1]$ — the second kind of
damping, used when the Taylor model is unreliable far from the optimum. It keeps Newton's
curvature-aware direction but covers only a fraction of the recommended distance, so the single clean
jump is replaced by a controlled ``staircase'' of partial steps that walk towards the optimum and
re-Taylor as they go.

\begin{trapbox}[``the Hessian is positive definite'' — only proved for linear regression]
The PD proof applies to the \emph{specific} loss $\|Xw-y\|^2+\lambda\|w\|^2$ ($H=2(X^\top X+\lambda I)$,
PD everywhere). Newton's \emph{method} applies to any $C^2$ $f$. On a non-convex surface $H$ can be
\kw{indefinite} (mixed-sign eigenvalues) — still \kw{invertible} if no eigenvalue is exactly zero,
but not PD. E.g. $\big(\begin{smallmatrix}1&0\\0&-1\end{smallmatrix}\big)$: $\det=-1$, invertible,
not PD. So invertible $\ne$ convex.
\end{trapbox}

\subsection{Why GD usually beats Newton in practice}

\begin{center}
\renewcommand{\arraystretch}{1.2}
\begin{tabularx}{\linewidth}{l l X}
\toprule
& \thead{Gradient descent} & \thead{Newton} \\
\midrule
Per-iteration cost & $O(NM)$ & $O(NM^2+M^3)$ (Hessian build + invert) \\
Memory & $O(M)$ gradient & $O(M^2)$ Hessian ($\sim$100$\times$ for large $M$) \\
Convergence rate & linear & quadratic (but only \emph{local}) \\
\bottomrule
\end{tabularx}
\end{center}
For large $M$, Newton's $M^2$ memory and $M^3$ inversion are prohibitive — a single transformer layer
has $M\sim10^8$, so the Hessian alone would have $10^{16}$ entries, physically uncomputable. Hence
GD / SGD dominate at scale, and Newton-style methods survive there only as cheap \emph{approximations}
of $H^{-1}$. For unit-scale linear regression ($M\le1000$) full Newton is feasible and elegant; the
storage cost matters most for embedded deployment where running memory is tightly constrained.

\begin{center}\large\bfseries\color{defBd} PART III — Probabilistic Modelling (Weeks 10–11)\end{center}

% =====================================================================
\section{Week 10 — Bayesian Linear Regression}

Week 10 is not a new problem — it is Week 8's linear regression \emph{re-told as a story about
probability}. The data $X$ and labels $y$ are the same; what changes is that we now imagine the
labels were \emph{generated} by a noisy process and ask which parameters were most likely to have
produced them. This reframing is worth the effort because it pays off twice over. First, ordinary
least squares \kw{falls out for free}: under Gaussian noise, the maximum-likelihood estimator turns
out to be exactly Week 8's squared-error minimiser, so the closed-form solution you already know is
revealed to be a principled estimator rather than an arbitrary heuristic. Second, regularisation
\kw{falls out for free}: bolting a prior onto the parameters and doing MAP estimation reproduces
ridge regression exactly — which is the bridge to Week 21. The deeper, fully Bayesian extension
(predictive distributions, uncertainty) is flagged non-examinable, so these notes stop at the two
point estimators.

\begin{setupbox}[probabilistic model]
$y_i=w^\top x_i+\epsilon_i$ with \kw{i.i.d.\ Gaussian noise} $\epsilon_i\sim\mathcal N(0,\sigma^2)$
(independent, identically distributed; zero mean = unbiased; $\sigma^2$ = spread). \\
Equivalently $\;y_i\mid x_i,w\sim\mathcal N(w^\top x_i,\sigma^2)$ — each label = the line plus noise.
\end{setupbox}

The story behind this model: there is some true function producing the labels, but every real
measurement deviates from it for a thousand small reasons — sensor jitter, rounding, stray
environmental effects. We bundle all of them into a single random variable $\epsilon$ and assume the
truth itself is linear. ``I.i.d.'' means the noise on each sample is drawn independently and from the
same distribution; the Gaussian is the standard default. Its zero mean encodes that the noise does
not systematically push observations up or down, and $\sigma^2$ sets the width of the bell: as
$\sigma\to0$ the distribution collapses onto the mean (every label is exactly the line) and as
$\sigma$ grows it flattens towards pure randomness. Crucially, once $x_i$ and $w$ are fixed the
\emph{only} randomness left in $y_i$ is the noise, so subtracting the prediction gives
$y_i-w^\top x_i=\epsilon_i\sim\mathcal N(0,\sigma^2)$ — which is exactly the equivalent form in the
box: each label is a Gaussian centred on the model's prediction.

\subsection{Likelihood and Maximum Likelihood}

\begin{defbox}[likelihood]
$p(y\mid X,w)$ = probability of the \kw{observed data}, treating a parameter $w$ as the truth.
A \kw{function of $w$}, with data $X,y$ \kw{fixed}.
\end{defbox}

The likelihood is the object you slide $w$ around to maximise: the data is nailed down (it is what
you actually collected) and $w$ is the free variable. This is the single most-misread idea in the
topic, so it gets a trap of its own.

\begin{trapbox}[likelihood is not P(parameters correct) — it is the reverse]
Likelihood = P(seeing this exact data $\mid$ $w$ assumed true), \emph{not} P($w$ is right).
\kw{Coin}: 9 heads in 10 flips is the fixed data; likelihood of ``fair'' ($w{=}0.5$) is low,
likelihood of ``weighted'' ($w{=}0.9$) is high. \kw{ML} = pick the $w$ that makes the observed data
least surprising.
\end{trapbox}

Read the coin example carefully, because it pins down the direction of the conditional. The 9-heads
run is the fixed evidence. You then score \emph{candidate} coins by it: a fair coin would rarely
produce so lopsided a run, so its likelihood is low; a 90\%-heads coin produces exactly such runs,
so its likelihood is high. At no point are you computing ``the probability the coin is fair'' —
that would require a prior. Maximum likelihood simply returns the parameter under which the observed
data is least surprising. Note too that the likelihood is a function of $w$ and need \emph{not}
integrate to 1 over $w$; it is not a distribution over parameters.

Because the samples are i.i.d., the probability of the whole dataset factorises into a product of
per-sample Gaussians, and that product collapses neatly because the exponents add:
\[
p(y\mid X,w)=\prod_{i=1}^N p(y_i\mid x_i,w)
=\frac{1}{(2\pi\sigma^2)^{N/2}}\exp\!\Big(-\frac{\|Xw-y\|^2}{2\sigma^2}\Big).
\]
The numerator inside the exponential, $\sum_i(y_i-w^\top x_i)^2$, is precisely the sum of squared
errors $\|Xw-y\|^2$ from Week 8 — that is the first clue that least squares is hiding inside this
Gaussian. \kw{Maximum Likelihood} then picks $w_\ast=\arg\max_w p(y\mid X,w)$. The clean way to do
this is to take the log first: the logarithm is monotone (so the $\arg\max$ is unchanged), it turns
the product into a sum, and it kills the exponential, leaving
$\log p=-\tfrac{N}{2}\log(2\pi\sigma^2)-\tfrac{1}{2\sigma^2}\|Xw-y\|^2$. The first term does not
depend on $w$ and drops out; maximising the second is the same as minimising the squared error:
\[
w_\ast=\arg\min_w\|Xw-y\|^2 \;=\;(X^\top X)^{-1}X^\top y.
\]
\examflag{}\kw{``Under i.i.d.\ Gaussian noise, ML estimation $\equiv$ ordinary least squares.''}
Wk8's closed form \emph{is} the ML estimator — not a heuristic. This is the headline result of the
lecture's first half and a likely scenario-question target.

\subsection{Bayes' rule, prior, posterior, MAP}

\kw{Why move from ML to MAP}: ML needs a lot of data, and on a small or unrepresentative sample it
overfits. The lecturer's pointed example: flip a coin three times, see three heads, and ML concludes
the coin \emph{always} lands heads — because that is literally the parameter under which the data is
most probable. ML has no way to push back against a freak sample. A prior is exactly that pushback:
it injects external belief about $w$ that regularises the freak conclusion away. MAP estimation is
the mechanism for folding that belief in, via Bayes' rule.

\begin{defbox}[Bayes' rule for parameters]
\[
\underbrace{p(w\mid X,y)}_{\text{posterior}}=
\frac{\overbrace{p(y\mid X,w)}^{\text{likelihood}}\;\overbrace{p(w)}^{\text{prior}}}
{\underbrace{p(y\mid X)}_{\text{evidence}}}
\;\;\propto\;\; p(y\mid X,w)\,p(w).
\]
\kw{Prior} = belief about $w$ \emph{before} data. \kw{Posterior} = updated belief \emph{after} data.
\kw{Evidence} = normaliser, constant in $w$, dropped for $\arg\max$.
\end{defbox}

The point of the proportionality is practical: the evidence $p(y\mid X)$ is a constant in $w$ (it is
just the normaliser that makes the posterior integrate to 1), so for the purpose of finding the
\emph{location} of the posterior's peak it can be ignored entirely. The prior encodes what you
believe about $w$ before seeing any data; the posterior is that belief after the data has spoken;
the likelihood is the data's verdict. Bayes' rule is simply the rule that reconciles the two.

\begin{defbox}[MAP estimation]
$w_\ast=\arg\max_w p(w\mid X,y)=\arg\max_w p(y\mid X,w)\,p(w)$ — most probable $w$ \emph{after}
seeing the data.
\end{defbox}

So MAP is just ML with one extra factor — the prior — multiplied into the objective. The likelihood
pulls $w$ towards parameters that explain the data; the prior anchors $w$ towards what you assumed
beforehand; MAP picks the single peak where those two pulls balance.

\begin{center}
\renewcommand{\arraystretch}{1.2}
\begin{tabularx}{\linewidth}{l X X}
\toprule
& \thead{ML} & \thead{MAP} \\
\midrule
Maximises & likelihood $p(y\mid X,w)$ & posterior $p(w\mid X,y)$ \\
Uses a prior? & no & yes \\
Reduces to & least squares & ridge / $L_2$-regularised least squares \\
\bottomrule
\end{tabularx}
\end{center}

MAP $=$ ML when the prior is \kw{uniform} (uninformative) — with a constant prior the $p(w)$ factor
contributes nothing to the $\arg\max$ and drops out, leaving the bare likelihood. So ML is not a
rival method but the special case of MAP in which you bring no prior information at all.

\subsection{MAP with a Gaussian prior $=$ ridge}

The simplest non-trivial prior is the zero-mean isotropic Gaussian $p(w)=\mathcal N(0,I)$, which in
plain English says ``before seeing the data I expect the weights to be small, and I am equally
suspicious of any one component being large in any direction.'' It is the natural default when you
have no domain knowledge about $w$. Substituting the Gaussian likelihood and the Gaussian prior into
the posterior, taking the log, and dropping the terms constant in $w$ gives
$\log p(w\mid X,y)=-\tfrac{1}{2\sigma^2}\|Xw-y\|^2-\tfrac12\|w\|^2+\text{const}$ — a sum of two
terms, the data-fit term and the prior penalty. Multiplying through by $-2\sigma^2$ (which flips the
$\arg\max$ to an $\arg\min$ and relocates the $\sigma^2$ onto the penalty) gives
\[
w_\ast=\arg\min_w\big(\|Xw-y\|^2+\sigma^2\|w\|^2\big)\;=\;(X^\top X+\sigma^2 I)^{-1}X^\top y.
\]
The second equality is just Week 8's differentiate-and-set-to-zero applied to the new objective:
$\nabla=2X^\top Xw-2X^\top y+2\sigma^2 w=0$ rearranges to $(X^\top X+\sigma^2 I)w_\ast=X^\top y$.
\examflag{}\kw{``MAP with a Gaussian prior $\equiv$ ridge / $L_2$-regularised least squares, with
regularisation strength $\sigma^2$ (the noise variance).''} Forward link to Wk21.

\begin{notebox}[why $\sigma^2$ appears in MAP but cancels in ML]
\kw{ML}: only goal is to fit the data; with no prior to balance against, $\sigma^2$ is just the
\emph{unit} the error is measured in $\Rightarrow$ it scales the objective uniformly and
\kw{cancels out of the $\arg\min$}. \\
\kw{MAP}: a \kw{tug-of-war} — likelihood pulls $w$ to fit the data, prior pulls $w$ toward $0$.
$\sigma^2$ is the \kw{trust factor} setting who wins. Large $\sigma^2$ (noisy data) $\Rightarrow$
stronger penalty $\Rightarrow$ lean on the prior. Small $\sigma^2$ (clean data) $\Rightarrow$ penalty
vanishes $\Rightarrow$ behaves like OLS.
\end{notebox}

\textbf{What the prior buys:} (i) \kw{numerical stability} — $X^\top X+\sigma^2 I$ is \emph{always}
invertible for $\sigma^2>0$ (every eigenvalue lifted by $\sigma^2$), even when $X^\top X$ is singular;
(ii) \kw{bias toward small weights} $\Rightarrow$ smoother fit, less overfitting.

\begin{pitfallbox}[ML and MAP are point estimates — no uncertainty]
Both return a single best $w_\ast$ $\Rightarrow$ a single-number prediction, \kw{no margin of error}.
\kw{Full Bayesian inference} keeps the whole posterior over $w$ and integrates it out to give a
\kw{predictive distribution} (``most likely 5, plausibly 3–7''). Optional / non-examinable.
\end{pitfallbox}

\begin{exambox}[Wk10 — the two derivations]
\textbf{ML $\to$ least squares (4 marks):} model $y_i\mid x_i,w\sim\mathcal N(w^\top x_i,\sigma^2)$;
likelihood $\propto\exp(-\|Xw-y\|^2/2\sigma^2)$; log-likelihood $=-\tfrac{1}{2\sigma^2}\|Xw-y\|^2+$const;
drop constants, flip sign $\Rightarrow\arg\min\|Xw-y\|^2$. \\
\textbf{MAP $\to$ ridge (5 marks):} posterior $\propto$ likelihood $\times$ prior $\mathcal N(0,I)$;
log-posterior $=-\tfrac{1}{2\sigma^2}\|Xw-y\|^2-\tfrac12\|w\|^2+$const; $\times(-2\sigma^2)$
$\Rightarrow\arg\min(\|Xw-y\|^2+\sigma^2\|w\|^2)$; differentiate, set $0$
$\Rightarrow w_\ast=(X^\top X+\sigma^2 I)^{-1}X^\top y$.
\end{exambox}

% =====================================================================
\section{Week 11 — Logistic Regression}

Logistic regression is the unit's first dedicated classifier built on the linear-model machinery.
Decision trees and random forests classified by recursively partitioning the input space; logistic
regression instead fits a smooth probability function over it. The trick is to take the familiar
linear score $w^\top x$ and squash it through a sigmoid so the output becomes a probability in
$(0,1)$, which can then be thresholded into a class. Everything else — the i.i.d.\ likelihood, the
$\arg\max$ on the log-likelihood, the gradient-ascent loop — is recycled from Weeks 8 and 10; Week 11
adds only the squashing step, the binomial likelihood, and the resulting linear decision boundary.

\begin{setupbox}[binary classification]
Data $\{(x_i,y_i)\}$, labels $y_i\in\{0,1\}$. Goal: predict $\hat y\in\{0,1\}$ for a new $x$.
Logistic regression squashes the linear score $w^\top x$ through a sigmoid to a probability.
Despite the name it is a \kw{classifier} — ``regression'' refers to regressing the \emph{log-odds}
on the features. Formulae below assume the $\{0,1\}$ convention; a $\{-1,+1\}$ label convention
changes them — check first.
\end{setupbox}

\subsection{Why not linear regression}

The natural first thought is to fit a plain linear model $f(x)=w^\top x$ to the $\{0,1\}$ labels and
threshold at $0.5$. It fails for two structural reasons, and the obvious patches fail too.

\begin{itemize}
  \item \kw{Unbounded output}: $w^\top x$ returns any real number — for an extreme input the
  prediction can be $-3$ or $+12$, which is meaningless as a probability or a class label. The model
  has no notion that the output should be confined.
  \item \kw{Outlier tilt}: squared loss penalises an extreme correct point (predicts $5$ for label
  $1$ $\Rightarrow$ loss $(5-1)^2=16$) $\Rightarrow$ the line tilts flatter to reduce that fake error
  $\Rightarrow$ the $0.5$-crossing decision boundary shifts $\Rightarrow$ normal middle points
  misclassified. The subtlety worth holding onto: the point was classified \emph{correctly}
  ($5>0.5$), yet squared loss still charges a huge error for it, and relieving that fake error
  sabotages the threshold for everyone else. The classifier is wrecked by a point it got right.
  \item Hacks fail: \kw{clipping} to $[0,1]$ discards variance info — a prediction of $5$ and one of
  $1.01$ both flatten to $1$, losing all sense of confidence, and the clipped regions have zero
  gradient; \kw{hard thresholding} at $0.5$ makes the function a discontinuous step
  $\Rightarrow$ non-differentiable at the jump and flat elsewhere, so gradient descent has no slope
  to follow.
  \item \kw{Fix}: compose $w^\top x$ with a smooth monotone squashing function that maps
  $\mathbb R\to(0,1)$ — the sigmoid. It saturates gently towards $0$ and $1$ so far-out points stop
  generating runaway error, yet it stays differentiable everywhere so training still works. The
  result is logistic regression.
\end{itemize}

\subsection{Sigmoid and logit}

\begin{defbox}[sigmoid / logistic function]
$\sigma(z)=\dfrac{1}{1+e^{-z}}$, maps $\mathbb R\to(0,1)$. Monotone, smooth ($C^\infty$),
$\sigma(0)=0.5$, $\sigma(z)+\sigma(-z)=1$. \\
\kw{Self-derivative}: $\sigma'(z)=\sigma(z)(1-\sigma(z))$ — memorise; used in every gradient.
\end{defbox}

The sigmoid's six properties are not trivia — each earns its place. The range $(0,1)$ makes the
output a valid probability; monotonicity preserves the order of inputs as an order of
probabilities; $\sigma(0)=0.5$ means the decision boundary $\sigma(w^\top x)=0.5$ is the clean
hyperplane $w^\top x=0$; smoothness ($C^\infty$) is what lets gradient-based training work;
the symmetry $\sigma(z)+\sigma(-z)=1$ says class-0 probability mirrors class-1; and the
self-derivative is the algebraic gift that makes every gradient collapse to something clean.

\begin{defbox}[logit / log-odds — inverse sigmoid]
$\text{logit}(p)=\log\dfrac{p}{1-p}$, maps $(0,1)\to\mathbb R$. \kw{Odds} $=p/(1-p)$; logit = log of
the odds. The model asserts $\;\log\dfrac{P(y{=}1\mid x)}{P(y{=}0\mid x)}=w^\top x$.
\end{defbox}

The logit is the sigmoid run backwards: where the sigmoid squashes a real score into a probability,
the logit recovers the score from the probability. Building it up in two steps makes its meaning
clear — the odds $p/(1-p)$ rescale a probability to $(0,\infty)$ (a $p=0.75$ becomes odds of 3,
``3 to 1 on class 1''), and the logarithm then stretches $(0,\infty)$ out to the whole real line.
So the assertion ``the log-odds of class 1 are linear in the features'' is the actual modelling
claim logistic regression makes.

\begin{exbox}[deriving $\sigma'$ and the logit]
\kw{Self-derivative}: $\sigma(z)=(1+e^{-z})^{-1}$, so by the chain rule
$\sigma'(z)=\dfrac{e^{-z}}{(1+e^{-z})^2}=\dfrac{1}{1+e^{-z}}\cdot\dfrac{e^{-z}}{1+e^{-z}}
=\sigma(z)\,(1-\sigma(z))$ — since $\dfrac{e^{-z}}{1+e^{-z}}=1-\sigma(z)$.
\kw{Logit} = invert $p=\sigma(x)$: $p(1+e^{-x})=1\Rightarrow e^{-x}=\tfrac{1-p}{p}
\Rightarrow e^{x}=\tfrac{p}{1-p}\Rightarrow x=\log\tfrac{p}{1-p}$.
\end{exbox}

\begin{notebox}[what the real-number score $z=w^\top x$ means]
$z$ is the model's confidence in \kw{log-odds} units. $z=0\Rightarrow$ odds $1\Rightarrow p=0.5$ =
the \kw{decision boundary}. $z>0\Rightarrow p>0.5$ (class 1), larger $z$ = more confident.
$z<0\Rightarrow p<0.5$ (class 0). Each weight $w_j$: a 1-unit rise in $x_j$ raises the log-odds of
class 1 by $w_j$ — linear interpretability.
\end{notebox}

\subsection{Model, likelihood, training}

The model itself is just the linear score fed through the sigmoid; everything below is the standard
maximum-likelihood pipeline applied to that model, with one or two pieces of pleasant algebra.

\begin{itemize}
  \item \kw{Model}: $f_w(x)=\sigma(w^\top x)$; $P(y{=}1\mid x)=f_w(x)$, $P(y{=}0\mid x)=1-f_w(x)$.
  Conventionally $x$ is augmented with a leading $1$ so that $w_0$ acts as a bias.
  \item \kw{Compact likelihood} (exponent trick): $p(y\mid x;w)=[f_w(x)]^y[1-f_w(x)]^{1-y}$ — the
  exponent picks the right factor for $y\in\{0,1\}$. When $y=1$ the second factor's exponent is $0$
  so that factor becomes $1$, leaving $f_w(x)$; when $y=0$ the first factor vanishes, leaving
  $1-f_w(x)$. The exponents act as on/off switches, and writing it this way is what avoids an ugly
  case-split inside the log-likelihood.
  \item \kw{Log-likelihood}: $\log L(w)=\sum_i[y_i\log f_w(x_i)+(1-y_i)\log(1-f_w(x_i))]$. Take the
  i.i.d.\ product of compact likelihoods and log it — the product becomes a sum. This is the same
  object as the \kw{negative cross-entropy} loss used in deep learning, only with the sign flipped.
  \item \kw{Gradient}: $\nabla\log L(w)=X^\top(y-\boldsymbol\sigma)$ — feature matrix times prediction
  errors; same structural form as linear regression's gradient. \emph{Why}: differentiating one term
  needs $\partial f/\partial w_j=\sigma'(z)x_{ij}=f(1-f)x_{ij}$; the $1/f$ from $\log f$ then cancels
  the $f$, and the $-1/(1-f)$ from $\log(1-f)$ cancels the $(1-f)$, leaving the clean residual
  $\partial_{w_j}\log L=\sum_i(y_i-f_w(x_i))\,x_{ij}$. The gradient is the input matrix times the
  vector of prediction errors — read it as feature-weighted residuals, exactly as in linear
  regression, only with sigmoid residuals.
  \item \kw{No closed form} ($\sigma$ is transcendental in $w$, so $\nabla\log L=0$ is not a linear
  system) $\Rightarrow$ train iteratively: gradient ascent $w_{t+1}=w_t+\alpha\nabla\log L$, or
  Newton/IRLS (few iterations, $H=-X^\top S X$, $S=\text{diag}(\sigma_i(1-\sigma_i))$). IRLS —
  iteratively reweighted least squares — is what most statistics software actually runs.
  \item \kw{Concave} log-likelihood (Hessian negative semi-definite) $\Rightarrow$ every local
  maximum is global, so gradient ascent cannot get trapped the way it can on a non-convex loss.
  Caveat: concavity gives a global max but not a \emph{unique} one — uniqueness needs strict
  concavity, i.e.\ $X$ of full column rank.
\end{itemize}

\subsection{Decision boundary and threshold}

The natural decision rule is to predict class 1 whenever $f_w(x)\ge0.5$. By the sigmoid's symmetry
$f_w(x)=0.5$ exactly when $w^\top x=0$, so the \kw{decision boundary} is the hyperplane $w^\top x=0$
(a line in 2-D, a plane in 3-D, an $(M{-}1)$-dimensional hyperplane in general) — \kw{linear in input
space}, the same shape as the SVM boundary you meet in Weeks 25--26. Two things follow that students
routinely confuse: the bias $w_0$ \emph{shifts} where that hyperplane sits, whereas the slope (the
magnitude of $w$) only controls how sharply confidence changes \emph{across} the boundary — it does
\kw{not} relocate the boundary itself.

\begin{notebox}[weight tuning under class imbalance — still fitting $\sigma(w^\top x)$]
With imbalanced data, MLE \emph{learns} a shifted/shallower curve (large bias, smaller slope) to give
the minority class fair probability space. The model form never changes — it is always
$\sigma(w^\top x)$; ``tuning'' = gradient ascent naturally finding the best $w$ for skewed data.
\end{notebox}

\begin{notebox}[moving the $0.5$ threshold is a choice, not retraining]
Changing the cutoff (e.g.\ to $0.2$ or $0.9$) does not move the sigmoid or retrain — it only changes
how strictly confidence is acted on, trading FP against FN. \kw{Lower} (cancer screening: a false
negative is disastrous $\Rightarrow$ flag at $p\ge0.2$). \kw{Raise} (spam filter: a false positive
loses a real email $\Rightarrow$ flag only at $p\ge0.9$).
\end{notebox}

\subsection{What it buys you, and its limitations}

\kw{Advantages}: (i) \kw{probabilistic output} — a calibrated $P(y{=}1\mid x)$, not just a label, so
downstream decisions can act on a confidence rather than a bare yes/no; (ii) \kw{efficient training}
— the log-likelihood is concave, the gradient is well-behaved, and Newton/IRLS converges in a
handful of iterations; (iii) \kw{interpretability} — each weight $w_j$ is the log-odds effect of
feature $j$, so the model can be read and explained to a non-expert.

\kw{Limitations}:
\begin{itemize}
  \item \kw{Linear boundary only} — the boundary is always the hyperplane $w^\top x=0$, so data with
  genuinely non-linear structure (e.g.\ a class enclosed in a circular ring) cannot be perfectly
  separated by any $w$. Fix: feature engineering ($x\to[\dots,x_1^2,x_1x_2,\dots]$ makes the
  boundary linear in the \emph{new} feature space) or a richer model (RF, kernel SVM, NN).
  \item \kw{Overfitting on (near-)separable data}: on separable data some $w$ classifies every
  point correctly, but maximum likelihood is greedy for confidence — it does not just want
  $p>0.5$, it wants $p$ as close to $1$ as possible because that is what makes $\log L$ larger. To
  push the sigmoid output to exactly $1$ needs $z=w^\top x=+\infty$, and since the inputs are fixed
  the only lever is to scale $\|w\|$ up without bound. So the optimiser drives $\|w\|\to\infty$,
  \kw{no finite maximiser} exists, and the smooth S-curve hardens into a brittle near-vertical step.
  Fix: \kw{regularisation} $\arg\max(\log L(w)-\tfrac\lambda2\|w\|^2)$ — the penalty grows without
  bound as $\|w\|$ grows, restoring a finite trade-off point. This is exactly MAP with a Gaussian
  prior (Wk10).
\end{itemize}

\begin{trapbox}[binary only? — softmax is the multi-class generalisation]
Standard logistic regression is \kw{binary} (sigmoid gives one $p$, the other is $1-p$).
\kw{Multinomial logistic regression}: one linear score $z_k=w_k^\top x$ per class, combined by
\kw{softmax} $P(k)=e^{z_k}/\sum_j e^{z_j}$ — guarantees the class probabilities are positive and sum
to 1. Sigmoid = binary case of softmax.
\end{trapbox}

\subsection{Sensitivity to scaling and outliers}

\kw{Scaling} — features on different scales cause: (i) an \kw{elliptical loss surface}
$\Rightarrow$ gradient ascent zig-zags, slow; (ii) \kw{unfair L2 penalty} — the regulariser treats
all weights equally, so a small-scale feature's (necessarily large) weight is over-penalised.
\kw{Fix}: standardise features (mean 0, variance 1).

\kw{Outliers / noisy labels} — cross-entropy loss is $-\log p$; as $p\to0$ the penalty $\to\infty$.
A mislabelled extreme point gets a huge penalty $\Rightarrow$ the decision boundary shifts to
accommodate one bad point. \kw{Fix}: robust losses that cap the penalty for confidently-wrong points.

\begin{center}
\renewcommand{\arraystretch}{1.25}
\begin{tabularx}{\linewidth}{l X X}
\toprule
\thead{Model} & \thead{Scaling sensitivity} & \thead{Outlier sensitivity} \\
\midrule
Linear regression & yes — GD zig-zags; regulariser unfair $\Rightarrow$ standardise &
yes — squared error, one outlier drags the line \\
Logistic regression & yes — same two reasons $\Rightarrow$ standardise &
yes — $-\log p$ penalty unbounded, boundary shifts \\
Bayesian regression & yes — one prior on all weights only fair if features comparably scaled &
yes — Gaussian noise model; needs heavy-tailed (Student-$t$) noise to be robust \\
Decision trees / RF & \kw{no} — splits depend only on value \emph{order}, scale-invariant &
\kw{robust} — an outlier is isolated in its own leaf, rest of tree unaffected \\
\bottomrule
\end{tabularx}
\end{center}

\begin{exambox}[Wk11 — past-paper Q4 reflexes]
\textbf{Why not linear regression}: unbounded output; squared-loss outlier tilt; clipping/thresholding
brittle $\Rightarrow$ use sigmoid. \quad
\textbf{$\sigma'$}: chain rule on $(1+e^{-z})^{-1}$ $\Rightarrow\sigma(z)(1-\sigma(z))$. \quad
\textbf{logit}: invert $p=\sigma(x)$ $\Rightarrow x=\log(p/(1-p))$. \quad
\textbf{log-likelihood}: i.i.d.\ product of $[f]^y[1-f]^{1-y}$, take log. \quad
\textbf{overfitting}: regularisation (L2 / Gaussian-prior MAP) stops $\|w\|\to\infty$.
\end{exambox}

\begin{center}\large\bfseries\color{defBd} PART IV — Unsupervised Learning (Week 19)\end{center}

% =====================================================================
\section{Week 19 — Clustering \& K-means}

Everything in the unit so far has been \emph{supervised}: every data point came with a label, and
the job was to predict that label for unseen inputs. Week 19 drops the labels. Now the dataset is
just inputs, and the task is to discover structure that nobody told the algorithm to look for.

\begin{setupbox}[unsupervised learning]
Data $D=\{x_1,\dots,x_N\}$ — \kw{inputs only, no labels $y$}. Goal: discover structure. Main tasks:
clustering, dimensionality reduction, density estimation. \kw{Why it matters}: most real data is
unlabelled — labels are expensive, often noisy, and frequently unavailable — so a method that needs
no labels is widely applicable.
\end{setupbox}

The reason unsupervised learning matters in practice is sheer availability: labelling requires human
annotation, which is slow and expensive, so the overwhelming majority of real-world data arrives
without labels. Unsupervised methods are what you reach for when you have abundant data but no
targets, or when the goal is to \emph{discover} structure rather than predict a known quantity.
Clustering — grouping similar samples — is the focus of this week.

\begin{defbox}[cluster]
A group of samples \kw{similar} to each other and \kw{less similar} to other groups, by some
distance criterion. \kw{$L_2$ / Euclidean} = default; \kw{$L_1$ / Manhattan} = robust to outliers.
\end{defbox}

Notice the definition is only as precise as the distance function it rests on — ``similar'' has no
meaning until you choose a metric. Euclidean ($L_2$) is the default and pairs naturally with
K-means (see below); Manhattan ($L_1$) sums absolute rather than squared differences, so it does not
amplify extreme coordinates and is the robust choice when outliers are a concern. Clustering shows
up in image segmentation (each pixel a data point), market segmentation, data compression (replace
each point by its cluster centre), anonymisation, and missing-value imputation. Algorithms come in
four families: \kw{centroid-based} (K-means), \kw{hierarchical} (agglomerative/divisive),
\kw{density-based} (DBSCAN — optional), \kw{distribution-based} (GMM/EM — optional).

\subsection{K-means algorithm}

K-means answers the clustering problem with a strikingly simple alternating loop: guess $K$ group
centres, assign every point to its nearest centre, then move each centre to the average of the
points that chose it — and repeat until nothing moves.

\begin{defbox}[K-means]
\textbf{1. Init} — pick $K$ random data points as centroids $c_k$. \\
\textbf{2. Assign} — each point to its nearest centroid: $\text{ind}_i=\arg\min_k\|x_i-c_k\|^2$. \\
\textbf{3. Update} — move each centroid to the mean of its cluster: $c_k=\frac{1}{|C_k|}\sum_{x\in C_k}x$. \\
\textbf{4. Converge} — repeat 2–3 until centroids stop changing.
\end{defbox}

Read in plain English: the initial centroids are $K$ randomly chosen data points acting as ``team
captains''; the assignment step puts each point on the team of its nearest captain (the $\arg\min$
returns the \emph{index} of the winner, not the distance); the update step walks each captain to the
middle of its current team; and the loop stops when no captain needs to move.

\begin{notebox}[why $L_2$ ``pairs naturally'' with K-means]
K-means minimises the \kw{within-cluster sum of squares} $J=\sum_k\sum_{x\in C_k}\|x-c_k\|^2$. The
squared $L_2$ distance has a \kw{clean (linear) gradient}, and the point minimising the summed squared
$L_2$ distance of a group \kw{is exactly its mean} — which is what the update step computes. With
$L_1$ the minimiser is the \kw{median} instead (a different algorithm, \kw{K-medians}).
\end{notebox}

\begin{notebox}[why K-means converges]
$J$ is the ``score''. \kw{Assignment} step minimises $J$ over memberships (centroids fixed);
\kw{update} step minimises $J$ over centroids (memberships fixed) — so \kw{every step decreases $J$ or
leaves it unchanged}. $J$ is \kw{bounded below by 0}, and there are only finitely many ($K^N$)
possible assignments $\Rightarrow$ K-means terminates — at a \kw{local} minimum (not necessarily
global). Practical termination: max iteration count, or centroid-movement below a threshold.
\end{notebox}

After convergence, the input space has been carved into $K$ regions — a \kw{Voronoi partition} in
which each point belongs to whichever centroid is nearest, and the cell boundaries are the
perpendicular bisectors between centroid pairs. A new point is then classified in $O(K)$ distance
checks, making prediction cheap.

\begin{itemize}
  \item Result = a \kw{Voronoi partition} (each point to nearest centroid; boundaries = perpendicular
  bisectors). Complexity \kw{$O(KNM)$} per iteration — for each of $N$ points, a distance to each of
  $K$ centroids in $M$ dimensions. The number of iterations is typically small (tens to hundreds).
  \item \kw{Advantages}: guaranteed finite convergence; simple to implement (two alternating steps,
  very few hyperparameters); scales linearly in $N$; adaptive — a new point is assigned to its
  nearest centroid without re-running the whole algorithm.
  \item \examflag{}\kw{Disadvantages} (examinable): need $K$ in advance (no automatic mechanism — a
  wrong $K$ gives bad clusters); non-deterministic (the result depends on random init, so different
  seeds give different clusterings); sensitive to outliers (squared $L_2$ amplifies them and the
  mean update drags the centroid towards a far-flung point); poor on varying size/density (K-means
  implicitly assumes roughly spherical, equal-density clusters); scales badly in high dimension $M$
  (the curse of dimensionality makes all pairwise distances become similar and uninformative).
  \item \kw{Best-of-$N$}: run K-means $N$ times with different inits, keep the run with smallest
  total point-to-centroid distance — the standard fix for the local-minimum / bad-init problem. It
  is cheap precisely because one K-means run is fast, so even $N=100$ restarts stay affordable; this
  is what industry actually does.
\end{itemize}

\subsection{Hierarchical clustering}

Where K-means commits to a single flat partition into a fixed $K$, hierarchical clustering builds a
whole \kw{dendrogram} — a tree of nested clusters — which you then cut at any level to read off a
clustering of \emph{any} $K$, chosen \emph{after} you have seen the data rather than before.
\begin{itemize}
  \item \kw{Agglomerative} (bottom-up): start with $N$ singleton clusters and repeatedly merge the
  two closest. Cheaper, because it works greedily on local distances; it is the practical default.
  \item \kw{Divisive} (top-down): start with one cluster containing everything and repeatedly split
  the most heterogeneous (largest average internal dissimilarity). Expensive — it must weigh all
  possible splits — but each split decision uses a global view of the data.
  \item \kw{Linkage} = how cluster–cluster distance is measured: single (the minimum pairwise
  distance — causes \kw{chaining}), complete (the maximum pairwise distance — produces compact
  clusters), average (mean pairwise distance), centroid (distance between centroids).
  \item \kw{Chaining} (single linkage) — because single linkage merges on the closest \emph{pair}
  of members, a cluster can grow link-by-link with no control over its overall shape: each new
  point is absorbed for being near just \emph{one} member, so the cluster straggles outward into a
  chain and can even curl into a ring. Mitigate with semi-supervised guide points (a few known
  members per intended cluster), best-of-$N$ random starts, or manually-spread seed points (e.g.\
  convex-hull / polygon-corner seeding to cover the whole region). \kw{Trap}: the lecturer's term
  ``nearest-neighbour clustering'' means single-linkage clustering — it is \emph{not} KNN.
\end{itemize}

\subsection{Evaluating a clustering}

Without ground-truth labels, judging a clustering is itself a hard question. The lecture gives the
standard intuition and three metrics that formalise it.

\begin{defbox}[good clustering]
\kw{Low intra-cluster distance} (points within a cluster close) \kw{and} \kw{high inter-cluster
distance} (clusters well-separated) — both at once.
\end{defbox}

\begin{itemize}
  \item \kw{WCSS / inertia} $=\sum_k\sum_{x\in C_k}\|x-c_k\|^2$ — K-means' own objective; lower is
  better. But it \emph{always} falls as $K$ rises (reaching exactly $0$ when every point is its own
  cluster), so minimising WCSS alone would absurdly recommend $K=N$ — it cannot itself pick $K$.
  \item \kw{Silhouette score} (per point): $s=\dfrac{b-a}{\max(a,b)}$, where $a$ is the point's mean
  distance to others in its own cluster (how well it fits where it is) and $b$ its mean distance to
  the nearest \emph{other} cluster (how well it would fit the best alternative). Range $[-1,1]$;
  \kw{higher is better} — you want $b\gg a$, which pushes $s$ towards $+1$; $s\approx0$ means a
  boundary point and $s<0$ flags a likely mis-assignment.
  \item \kw{Davies–Bouldin index}: $DB=\dfrac{1}{K}\sum_i\max_{j\ne i}\dfrac{S_i+S_j}{M_{ij}}$,
  where $S_i$ is cluster $i$'s dispersion and $M_{ij}$ the separation of centroids $i$ and $j$.
  \kw{Lower is better} — you want tight clusters (small $S$ in the numerator) that are far apart
  (large $M$ in the denominator). Mnemonic: for each cluster take its \emph{worst} (most
  overlapping) neighbour, then average those worst-case ratios.
\end{itemize}

\begin{defbox}[elbow method — choosing $K$]
Plot WCSS vs $K=1,2,3,\dots$. The curve drops steeply then bends; the \kw{elbow} (the bend) marks a
good $K$. Beyond it, extra clusters give only diminishing WCSS reduction (just slicing real clusters).
\end{defbox}

The logic of the elbow: at small $K$ each new cluster captures a genuine natural grouping, so WCSS
falls sharply; once the real structure has been captured, extra clusters only slice already-tight
clusters and buy a tiny WCSS reduction each. The bend between those two regimes is the recommended
$K$ — it reads a sensible answer off the \emph{shape} of the curve, sidestepping the fact that the
WCSS \emph{value} would always prefer more clusters.

\begin{center}\large\bfseries\color{defBd} PART V — Parametric \& Non-parametric (Week 20)\end{center}

% =====================================================================
\section{Week 20 — Parametric \& Non-parametric Models}

Week 19's split — supervised vs unsupervised — asked whether you have labels. Week 20 introduces a
second, completely independent split that asks a different question: how does the model's complexity
scale with the amount of data? The two splits are \kw{orthogonal} — the same task can be tackled by
a parametric or a non-parametric method, supervised or unsupervised — so they form a 2×2 grid rather
than a single ranking.

\begin{defbox}[parametric vs non-parametric]
\kw{Parametric} — commits to a fixed functional form (linear, Gaussian, sigmoid\,\dots) with a
\kw{fixed, finite} number of parameters; \#params does \emph{not} grow with $N$.
e.g.\ linear/Bayesian/logistic regression, MLP, Laplace approximation. \\
\kw{Non-parametric} — no/weak distributional assumption; model complexity (effective \#params)
\kw{grows with $N$}. e.g.\ KDE, K-nearest-neighbours, kernel SVM, Gaussian processes.
\end{defbox}

The distinction is about \emph{commitment}. A parametric model decides on a functional shape before
it sees the data — linear regression commits to $w^\top x$, with $M$ parameters whatever $N$ is —
and training only fits the parameters of that shape. A non-parametric model makes no such
commitment: feeding it more data lets it become a genuinely richer function, because its effective
parameter count climbs with $N$.

\begin{trapbox}[``non-parametric'' is a misleading name]
It does \emph{not} mean ``no parameters'' — these models often have many. It means the parameter
count is \kw{not bounded in advance}: it scales with the training set.
\end{trapbox}

\begin{itemize}
  \item \kw{Parametric}: + faster (training is parameter estimation — a closed form or a few
  iterations), + a small training set is enough (the functional form does the heavy lifting via its
  inductive bias), + robust to overfit (a low-capacity form cannot memorise idiosyncrasies anyway);
  $-$ constrained to the assumed form (if the truth is not linear, no $w$ can fix that), $-$ leans
  on strong distributional assumptions whose correctness it cannot check.
  \item \kw{Non-parametric}: + fits many functional forms, including ones you could not have
  specified in advance, + makes no/weak assumptions, + often higher performance on hard problems
  with abundant data; $-$ data-hungry (without prior structure you need many samples), $-$ slow
  (often $O(N^2)$ or $O(N^3)$), $-$ higher capacity means higher overfitting risk.
  \item \examflag{}Past-paper MCQ: non-parametric \emph{is} flexible (TRUE — the headline benefit)
  and \emph{is} overfitting-prone (TRUE — the headline cost); it is \emph{not} fast to converge
  (FALSE — the model grows with the data) and does \emph{not} need small datasets (FALSE — the
  opposite).
\end{itemize}

\begin{trapbox}[parametric $\ne$ simple; the hybrid middle]
``Parametric'' is not ``low-capacity'': a high-degree polynomial has a fixed parameter count yet
overfits wildly. Random forests (fixed-depth trees, but the forest grows with data) and deep nets
(fixed architecture, millions of params) sit in the \kw{hybrid middle} — for an MCQ ``which is
non-parametric'', answer from the canonical list (KDE, KNN, kernel SVM, GP).
\end{trapbox}

\subsection{K-Nearest Neighbours}

\begin{defbox}[KNN]
Supervised. A \kw{lazy learner}: training just memorises the dataset. Predict for new $x$: compute
distance to every stored point, take the $k$ nearest, then \kw{majority vote} (classification) or
\kw{average} (regression). Non-parametric — the ``model'' is the training set, so its size grows
with $N$.
\end{defbox}

KNN is the cleanest illustration of what ``non-parametric'' means: it never compresses the data into
a fixed equation at all. There is no training computation — the algorithm just stores everything —
and a prediction directly consults the stored points. With $N$ training points a prediction
references $N$ points, so the model's size grows exactly with the data, which is the defining
property.

\begin{trapbox}[$k$ is a hyperparameter; KNN $\ne$ K-means]
\kw{Parameters} are learned from data (e.g.\ regression weights); \kw{hyperparameters} are set by you
beforehand ($k$ in KNN, bandwidth $h$ in KDE). The parametric/non-parametric label depends only on
\emph{parameters} — so ``non-parametric'' still has settings to choose.
\kw{KNN} (supervised; $k$ = neighbours to vote) is \emph{not} \kw{K-means} (Wk19; unsupervised;
$K$ = clusters, fixed in advance).
\end{trapbox}

\subsection{Kernel Density Estimation}

KDE is the canonical non-parametric \emph{density} estimator: given samples from an unknown
distribution, estimate the density at any test point. The starting point is the humble histogram —
bin the space, count samples per bin, divide by $N\cdot h$ to turn counts into a density. That works
but is rough: the estimate is piecewise constant, jumps at bin edges, and depends sensitively on
where the bins fall. KDE is the smooth generalisation that cures all three.

\begin{defbox}[KDE]
Density at $x$ is $f(x)=\lim_{h\to0}\tfrac{F(x+h)-F(x-h)}{2h}$ (derivative of the CDF). With finite
data and fixed $h$ this becomes a histogram (count in $[x{-}h,x{+}h]$ over $2hN$). Replace the hard
in/out count with a smooth \kw{kernel} $G$ weighting by distance:
\[
\bar f_h(x)=\frac{1}{Nh}\sum_{i=1}^N G\!\Big(\frac{x-x_i}{h}\Big).
\]
A smooth bump at every sample, summed. Non-parametric: $N$ terms, grows with data.
\end{defbox}

The key move is to replace the histogram's hard ``in or out of $[x-h,x+h]$'' indicator with a smooth
weight that decays with distance from $x$ — that weight is the \kw{kernel} $G$. Each sample then
contributes a smooth bump centred on itself rather than a hard count, and the density estimate at
$x$ is the sum of all those bumps. Common kernels are the Gaussian (the standard choice), the
uniform/box kernel (which recovers the histogram), and the Epanechnikov kernel (mean-squared-error
optimal). Because the estimate has one term per sample, adding data adds capacity — the training set
literally is the model.

\begin{exbox}[KDE — worked example, box kernel]
Samples $\{-2.1,-1.3,-0.4,1.9,5.1,6.2\}$, $N=6$; uniform box kernel ($G=\tfrac12$ on $[-1,1]$, so
$\bar f_h(x)=\frac{\text{count}}{2hN}$ where count $=\#\{x_i:|x-x_i|\le h\}$).
\kw{At $x=0$, $h=2$}: $-1.3,-0.4,1.9$ lie within distance $2$ — count $=3$, so
$\bar f_2(0)=\frac{3}{2\cdot2\cdot6}=0.125$. \kw{At $x=-1$, $h=1$}: the window $[-2,0]$ contains
$-1.3,-0.4$ — count $=2$, so $\bar f_1(-1)=\frac{2}{2\cdot1\cdot6}=\tfrac16\approx0.167$. A Gaussian
kernel replaces the hard in/out count with a smooth distance-weighted bump from \emph{every} sample.
\end{exbox}

\begin{notebox}[bandwidth $h$ — what it does and why it is a hyperparameter]
$h$ sets the \kw{width} of each bump (not the number of points — that is $N$). It cannot be fixed at
$1$: it must rescale the kernels to the data's units, and it is the \kw{bias-variance knob} (small
$h$ = spiky/overfit, large $h$ = over-smoothed/underfit). It appears twice: inside $G$ as a distance
scaler, in the $1/Nh$ prefactor to renormalise area to 1. KDE never \emph{learns} $h$ — it is a
\kw{hyperparameter}, set by \kw{Silverman's rule} $h\approx1.06\,\sigma\,N^{-1/5}$ ($\sigma$ = scale,
$N^{-1/5}$ shrinks $h$ as data grows) or by cross-validation.
\end{notebox}

\kw{Curse of dimensionality}: to keep the same point density as dimensions grow, the required sample
count scales as $k^M$ — $k$ points cover a line, $k^2$ a square, $k^3$ a cube, $k^M$ an $M$-D space
— so the requirement is \emph{exponential} in dimension $M$. A fixed-width kernel window is a
segment in 1-D, a small disc in 2-D, a small ball in 3-D; as $M$ climbs that window shrinks to a
negligible fraction of the space and captures few or zero neighbours, so KDE outputs near-zero
everywhere unless $N$ also grows exponentially. This is the central reason non-parametric methods
``need a lot of data,'' and why practical KDE is limited to small $M$.

\subsection{Laplace approximation}

\begin{defbox}[Laplace approximation]
The \kw{parametric exemplar}: approximate an intractable target $P^\ast(\theta)$ by a \kw{Gaussian}
$Q^\ast$ centred at its \kw{mode}, with width set by the curvature (Hessian) there. Parametric — a
Gaussian is fixed by a mean and a spread. \examflag{}Principle only; the maths is non-examinable.
\end{defbox}

The intuition: take an unwieldy target distribution you can evaluate but not handle analytically,
find its mode (the peak), and pretend the log-density is locally a downward parabola there — which
is exactly a Gaussian. It is the parametric exemplar because once you commit to ``Gaussian,'' the
whole approximation is pinned down by a fixed handful of numbers (a mean and a spread) that do not
grow with the data — the polar opposite of KDE.

\begin{itemize}
  \item \kw{Limitations} — all three stem from forcing one symmetric Gaussian onto the target:
  bad on \kw{multimodal} targets (a Gaussian has one peak, so it locks onto a single mode and
  ignores the rest of the mass); bad on \kw{heavy-tailed / skewed} targets (a Gaussian is perfectly
  symmetric with thin, fast-decaying tails, so it cannot bend to a lopsided density or cover a long
  tail); only as good as the \kw{local Hessian at the mode} — it is nearsighted, drawing the whole
  Gaussian from the curvature at one point and ignoring the global shape, so a sharp peak on an
  otherwise broad distribution yields a Gaussian that is far too narrow.
\end{itemize}

\begin{notebox}[KL divergence and Variational Inference]
\kw{KL divergence} = a number scoring the mismatch between two distributions ($0$ = identical,
larger = worse). \kw{Variational Inference (VI)} generalises Laplace: it picks $Q^\ast$ by
\kw{minimising the KL divergence} to $P^\ast$ — KL is the measuring stick, VI the search. Unlike
Laplace it fits the \emph{whole shape} (not just curvature at the mode) and is \kw{not restricted to
Gaussians}.
\end{notebox}

\begin{center}\large\bfseries\color{defBd} PART VI — Regularisation (Week 21)\end{center}

% =====================================================================
\section{Week 21 — Regularisation}

Week 10 already showed one route to regularisation — MAP with a Gaussian prior produces an $L_2$
penalty. Week 21 widens the picture: it gives the non-probabilistic motivations for regularising at
all, and adds the $L_1$ (lasso) and combined ($L_1{+}L_2$, elastic net) variants. The starting
intuition is the four-point regression: four data points lying almost on a line can be fitted with
\emph{zero} training error by a line, a sine curve, a step function, or a wild polynomial — squared
error cannot tell them apart. A human instantly picks the line, but the data alone does not justify
that choice. Regularisation is how we write that human preference for simplicity into the
optimisation as an equation.

\begin{defbox}[regularised loss]
$L(\theta)=\underbrace{L_{\text{data}}(\theta)}_{\text{data fit}}+\lambda\underbrace{R(\theta)}_{\text{penalty on }\theta}$.
$\lambda\ge0$ = regularisation strength, a \kw{hyperparameter} (tuned on validation, not learned).
$\lambda=0$ = unregularised; $\lambda\to\infty$ collapses $\theta$ to the regulariser's preference
(usually $0$).
\end{defbox}

The two terms pull against each other: $L_{\text{data}}$ rewards fitting the training data, while
$\lambda R(\theta)$ fines complexity, and $\lambda$ sets who wins. Note $\theta$ is found by the
optimisation but $\lambda$ is \emph{not} — it is held fixed during training and tuned externally on
a validation set. Different choices of $R$ encode different definitions of ``simple.''

\kw{Four reasons to regularise} — you should be able to name and one-sentence each: (1) avoid
\kw{overfitting} — without a penalty a high-capacity model memorises noise as if it were signal and
generalises poorly; (2) \kw{ill-posed problem} — when many parameter values fit the data equally
well, the optimiser can drift between them and never settle, so the penalty supplies a tiebreaker;
(3) \kw{human understanding} — a penalty that pushes toward sparse or interpretable parameters makes
the model readable; (4) \kw{easier optimisation} — a smooth convex penalty removes flat regions and
spurious local minima from the loss surface, so gradient methods converge more reliably. Often
several apply at once.

\begin{notebox}[Occam's razor]
``The simplest explanation is usually the correct one.'' The 4-point regression: a line, a sine, a
wild polynomial all fit 4 near-collinear points with \kw{zero training error} — the data alone
cannot choose. Occam's razor is the tiebreaker (pick the line); \kw{regularisation is Occam's razor
written as an equation} — the $\lambda R(\theta)$ term fines complexity. Overfitting = ``an
unjustifiably complex explanation''.
\end{notebox}

\subsection{Ridge ($L_2$) and Lasso ($L_1$)}

The two workhorse regularisers differ only in the norm they penalise — squared $L_2$ versus
absolute $L_1$ — but that single change of norm has consequences that run all the way through to
whether the solution is sparse and whether a closed form exists.

\begin{defbox}[ridge vs lasso]
\kw{Ridge / $L_2$}: $L=\|Xw-y\|^2+\lambda\|w\|^2$; closed form
$w_\ast=(X^\top X+\lambda I)^{-1}X^\top y$. \\
\kw{Lasso / $L_1$}: $L=\|Xw-y\|^2+\lambda\|w\|_1$; \kw{no closed form} ($|w|$ non-differentiable at
$0$) — solved iteratively (coordinate descent, LARS).
\end{defbox}

Ridge has a closed form because its penalty $\|w\|^2$ is smooth: differentiating the loss and
setting it to zero gives the linear system $(X^\top X+\lambda I)w_\ast=X^\top y$, the normal
equation with $\lambda I$ added to the Gram matrix. Lasso has none, because $|w|$ has a kink at the
origin where the gradient is undefined — so $\nabla L=0$ is not a clean equation and the problem
must be solved iteratively.

\kw{Why $L_2$ helps}: (i) \kw{numerical stability} — adding $\lambda I$ lifts every eigenvalue of
$X^\top X$ by $\lambda$, so $X^\top X+\lambda I$ is always invertible for $\lambda>0$ even when
$X^\top X$ is rank-deficient; (ii) reduces overfitting — penalising large $\|w\|$ means the model
cannot memorise idiosyncrasies that require large coefficients, so the fit is smoother; (iii)
\kw{shrinks but does not zero out} — each weight is pulled towards zero but never reaches it
exactly, so $L_2$ is \kw{not sparse}.

\begin{trapbox}[why $L_1$ is sparse and $L_2$ is not]
\examflag{}Past-paper: ``increasing $\lambda$ in lasso $\Rightarrow$ sparser coefficients'' — TRUE. \\
\kw{Geometric}: the optimum is where the data-fit ellipse meets the penalty ball. The $L_1$ ball is a
\kw{diamond} with corners \emph{on the axes} ($w_j=0$); as $\lambda$ grows the diamond shrinks and
the optimum lands on a corner $\Rightarrow$ exact zeros. The $L_2$ ball is a \kw{circle} (no corners)
— the optimum slides along it, never snapping to an axis. \\
\kw{Algebraic}: $L_1$'s slope is a \kw{constant $\lambda$} everywhere (even at $0$) — enough to
cancel a small data-fit gradient and pin a coordinate at $0$. $L_2$'s slope is $2w\to0$ as $w\to0$ —
the push fades, so weights shrink but never reach $0$. \\
$\Rightarrow$ Lasso does \kw{feature selection} (zeroed features are dropped); ridge keeps all.
\end{trapbox}

\subsection{Elastic net, sample size, interpretation}

\begin{itemize}
  \item \kw{Elastic net}: $L=\|Xw-y\|^2+\lambda(\gamma\|w\|_1+(1-\gamma)\|w\|^2)$ — two
  hyperparameters ($\lambda$ = overall strength, $\gamma$ = $L_1$/$L_2$ mix; $\gamma=0$ is pure
  ridge, $\gamma=1$ pure lasso). It blends the two because pure lasso has a quirk: faced with a
  group of correlated features it arbitrarily keeps one and zeros the rest, whereas the $L_2$
  component encourages correlated features to share similar coefficients. Elastic net thus combines
  $L_1$ sparsity-on-irrelevant-features with $L_2$ stability-on-correlated-features; $\gamma=0.5$ is
  a sensible default if you do not want to sweep it.
  \item \examflag{}\kw{$\lambda$ vs sample size}: $N\uparrow\Rightarrow\lambda\downarrow$. The
  data-fit term is a sum over $N$ samples while the penalty is fixed, so as $N$ grows the data-fit
  term increasingly dominates and a smaller $\lambda$ achieves the same balance. Equivalently: more
  data means less need for the prior bias the regulariser supplies — the data alone is informative
  enough to pin down a good $w$. (The marker also accepts a computational-stability framing.)
  \item \kw{Probabilistic view}: ridge $=$ MAP with a \kw{Gaussian} prior on $w$ (Wk10); lasso $=$
  MAP with a \kw{Laplace} prior. The same $\lambda R$ term is read as a negative log-prior — the
  non-probabilistic view treats $\lambda$ as a free knob, the probabilistic view as a prior
  parameter.
  \item $\lambda$ is swept on the \kw{validation} set (train fits $\theta$, validation tunes
  $\lambda$, test is touched once at the end; train/val/test discipline, §4). The validation curve
  is U-shaped — too small $\lambda$ overfits, too large underfits — so pick the bottom of the U.
  The \kw{bias/intercept $w_0$ is conventionally not penalised}, since shrinking it would pull the
  average prediction towards zero for no good reason.
\end{itemize}

\begin{notebox}[other forms of regularisation — and their limits]
\kw{Early stopping}: halt iterative training before the training error bottoms out — model
complexity is implicitly capped by how long you trained. Common in deep learning, but the lecturer
calls it a ``bad hack'': if you need it your explicit regularisation is too weak — though it may
still be the least-bad option. \kw{Model limits}: a restricted model class is itself regularisation
(logistic regression can only draw a linear boundary). \emph{Limitation}: the amount is fixed by the
class — rarely the right amount, with no $\lambda$ to tune. \kw{Data quantity}: more data needs less
regularisation; infinite data needs none. Models have a goldilocks range — too little data fails,
beyond a point extra data stops helping.
\end{notebox}

\begin{center}\large\bfseries\color{defBd} PART VII — Graphical Models \& Inference (Weeks 22–23)\end{center}

% =====================================================================
\section{Week 22 — Graphical Models}

Up to now every model has been a single conditional $P(y\mid x)$ — features in, label out, all
features treated symmetrically. Week 22 steps back and asks how to encode \emph{structure}: the
domain knowledge that some variables influence each other and others do not. A graphical model is
that knowledge drawn as a picture and read as a probability distribution. The running lecture
example is train derailment: rather than learn $P(\text{crash}\mid 6\text{ features})$ as one giant
table, you notice that three driver features cause the train's \emph{speed}, and speed plus three
hardware features cause the crash — so the joint factorises into two smaller, more learnable
conditionals. The graph is just the bookkeeping for that factorisation.

\begin{setupbox}[graphical model]
A graph encoding the structure of a joint distribution over random variables (RVs). Nodes = RVs;
edges encode \kw{conditional (in)dependence}. The structure is a \kw{modelling assumption} from
domain knowledge — a conditional independence is a property of \emph{the graph you drew}, not an
intrinsic fact (add an edge and the claim changes).
\end{setupbox}

The "property of the graph you drew" point is worth dwelling on. The derailment graph claims that,
once you know the speed, the speed \emph{limit} tells you nothing extra about the crash — speed
"screens off" the upstream driver features. But that claim is true only because you drew no direct
edge from the limit to the crash. Add that edge and the independence evaporates. Conditional
independence is therefore something you \emph{assert} by your choice of structure, not a fact you
discover; the graph is a hypothesis about how the data was generated.

\begin{notebox}[why graphical models are useful — the advantages]
\kw{Lower-dimensional sub-models}: each factor is a conditional with fewer inputs, so the curse of
dimensionality is cut — $P(\text{speed}\mid\text{3 features})$ is far easier to estimate than
$P(\text{crash}\mid\text{6 features})$. \kw{More effective use of data}: a moderate dataset trains
many small factors instead of one giant table. \kw{Missing-data \& latent-variable tolerance}:
inference still runs with some variables unobserved — even ones \emph{never} observed (latent),
which a standard supervised model cannot do. \kw{Performance}: baking in domain knowledge
``hard-removes'' impossibilities, tightening the model. \kw{Headline}: \emph{every} ML algorithm has
an implicit graphical model — some structural assumption about how the data was generated; graphical
models simply make it explicit. (A fully-connected graph — every feature into $y$ — encodes no
conditional independence and is therefore useless; the useful graphical models are the ones that
impose structure.)
\end{notebox}

\subsection{Three representations}

The same joint distribution can be drawn in three notations, each emphasising a different aspect of
the structure. They are not rival theories — they describe the same conditional-independence facts —
but the choice matters in practice because some operations (learning, inference, sampling) are
easier in one notation than another.

\begin{defbox}[BN / factor graph / Markov network]
\kw{Bayesian network} — directed acyclic graph; joint $=\prod_v P(v\mid\text{parents}(v))$. \\
\kw{Factor graph} — bipartite: circles (RVs) + squares (factors); joint $\propto\prod_f F_f(\text{its RVs})$. \\
\kw{Markov network} — undirected; joint $\propto\prod_{(i,j)}F_{ij}(x_i,x_j)\prod_i A_i(x_i)$
(pairwise + unary potentials).
\end{defbox}

A Bayesian network reads naturally as "cause $\to$ effect" — though strictly the arrow only encodes a
conditional dependency, and "arrow $=$ causation" is an intuition pump, not a theorem. Its great
virtue is that \emph{the graph is the factorisation}: one local conditional per node. A Markov
network drops all direction and is the right tool for symmetric, spatial relationships — adjacent
pixels in an image, where no pixel "causes" its neighbour. A factor graph is the most explicit of
the three: it draws the factor functions as their own square nodes, so multi-variable dependencies
are visible directly without needing arrows. Code typically uses factor graphs precisely because the
data structure is uniform.

\kw{Conversions}: BN$\to$FG and MN$\to$FG are \kw{always possible}; FG$\to$BN/MN is \kw{not}. The
asymmetry is because the factor graph is the most general of the three — collapsing a directed or
undirected graph into explicit factors never loses information, but the reverse may not have an
unambiguous answer.

\begin{trapbox}[why FG$\to$Markov network can fail — even though both are undirected]
(i) \kw{Pairwise-only}: a Markov network allows only 2-variable potentials; a factor graph can have
\kw{higher-order factors} ($G(a,c,d)$) — an irreducible 3-way factor cannot be rebuilt from edges.
(ii) \kw{Loss of explicitness}: one 3-way factor $F(a,b,c)$ and three pairwise $f(a,b)f(b,c)f(a,c)$
are different factor graphs but map to the \emph{same} Markov triangle — reversing is ambiguous.
\end{trapbox}

\subsection{Reading the graph}

\kw{Shaded} = observed, \kw{unshaded} = unobserved. Shading means only "do we have the value" — not
cause vs effect, not graph position. This is the subtle part: an observed effect at the bottom of
the graph can be shaded (you \emph{hear} the sirens) while an unobserved cause at the top is unshaded
(you cannot see the weather). That asymmetry is exactly what makes a graphical model powerful — it
can take observed effects and infer \emph{backwards} to unobserved causes, something a standard
feed-forward model cannot do. The topology never changes; only the shading does, and different modes
re-shade the same graph: \kw{training} (all data RVs shaded), \kw{testing} (inputs shaded, target
unshaded), \kw{missing data} (some inputs unshaded). A \kw{latent} variable is one that is
\emph{always} unobserved in every mode, and splits into two flavours: \kw{hidden} (a real quantity
we know exists but cannot measure — the train's true speed, with no track sensor) versus
\kw{hypothetical} (a construct the model posits that may have no measurable physical counterpart — a
topic in a topic model).

\subsection{Naive Bayes}

\begin{defbox}[Naive Bayes]
Assumes features are \kw{conditionally independent given the label}:
$P(x\mid y)=\prod_i P(x_i\mid y)$. Star structure ($y\to x_i$). Predict via Bayes' rule:
$P(y\mid x)\propto P(y)\prod_i P(x_i\mid y)$.
\end{defbox}

Naive Bayes flips ordinary classification on its head. A plain regression/classification model has
every feature pointing into $y$ — fully connected, no conditional independence, so the graphical
model says nothing useful. Naive Bayes instead makes $y$ the \emph{parent} and each $x_i$ a leaf, a
star with arrows $y\to x_i$, and reads it causally as "the class label generates the features." The
assumption that the features are independent once you know the label is strong and rarely literally
true, but it buys a tractable factorisation that works surprisingly well as a baseline (text
classification is the standard example).

\begin{notebox}[discriminative vs generative]
\kw{Discriminative} models learn $P(y\mid x)$ directly (logistic regression, NN, SVM).
\kw{Generative} models learn $P(x\mid y)$ and $P(y)$, then invert with Bayes' rule — Naive Bayes.
It flips the usual direction because $P(x_i\mid y)$ is just a count, sidestepping the curse of
dimensionality of estimating $P(y\mid x_1,\dots,x_n)$ jointly. The licence to go this way is that the
joint expands two ways — $P(y)P(x\mid y)=P(x)P(y\mid x)$ — so estimating the easy side
($P(x\mid y)$) recovers the side you actually want ($P(y\mid x)$).
\end{notebox}

\begin{notebox}[expert systems = a Bayesian network with Bernoulli factors]
An \kw{expert system} is a Bayesian network whose structure and probabilities are elicited from a
human expert rather than learned from data. If every variable is binary, every local conditional is
a \kw{Bernoulli}: with $\theta_{v\mid\text{pa}}=P(v{=}1\mid\text{parents})$, each factor is
$P(v\mid\text{pa})=\theta^{v}(1-\theta)^{1-v}$ (the exponent picks the right value for $v\in\{0,1\}$).
The full joint is the product over nodes, e.g.\
$P(s,c,b,x,d)=P(s)\,P(c\mid s)\,P(b\mid s)\,P(x\mid c)\,P(d\mid c,b)$ — each factor a Bernoulli with a
handful of expert-supplied parameters.
\end{notebox}

Learning Naive Bayes is therefore just counting, and predicting is just adding logs.

\begin{itemize}
  \item \kw{MLE}: $p_{iy}=C_i(1,y)/[C_i(0,y)+C_i(1,y)]$ — the empirical frequency of feature $i$
  being on within class $y$. It falls straight out of differentiating the Bernoulli log-likelihood
  $C_i(1,y)\log p_{iy}+C_i(0,y)\log(1-p_{iy})$ and setting the derivative to zero. The danger: the
  prediction is a \emph{product} of per-feature probabilities, so a single zero count anywhere
  zeroes the whole thing — a feature that never co-occurred with a class in training would veto that
  class entirely. Hence smoothing is needed.
  \item \kw{Test-time}: work in \kw{log space} —
  $\log P(y\mid x)=\log P(y)+\sum_i[x_i\log p_{iy}+(1-x_i)\log(1-p_{iy})]$. A long product of small
  probabilities underflows to zero in floating point; taking logs turns the product into a sum,
  which is numerically safe, and you simply pick the $y$ with the largest log-score.
\end{itemize}

\begin{trapbox}[Laplace smoothing = posterior \emph{mean} of a uniform Beta prior]
Beta is the \kw{conjugate prior} for Bernoulli: posterior $=\text{Beta}(\alpha+C_1,\beta+C_0)$ —
conjugacy meaning prior and posterior have the same functional form, so updating just adds the
counts to the shape parameters.
\kw{Laplace (add-one)} $p_{iy}=\frac{C_1+1}{C_0+C_1+2}$ is the posterior \kw{mean} with
$\alpha=\beta=1$. It is \kw{not the mode}: the mode with $\alpha=\beta=1$ cancels back to the
unsmoothed MLE — you would need $\alpha=\beta=2$. So Laplace smoothing is not an arbitrary hack; it
is exactly Bayesian estimation under a uniform prior, read off as the posterior mean. (The lecture
slide misstates this as the mode.)
\end{trapbox}

\subsection{Explaining away and d-separation}

The single most counter-intuitive pattern in graphical models — and a favourite of MCQs — is the
V-structure. The everyday rule of thumb is "two competing causes are independent," and marginally
that is true. The unfamiliar move is that \emph{conditioning on a common effect creates dependence}
— the opposite of the usual "conditioning makes things independent" reflex that holds for chains and
forks.

\begin{defbox}[V-structure / collider — explaining away]
$a\to z\leftarrow b$: parents $a,b$ are \kw{marginally independent} but become \kw{conditionally
dependent} once the child $z$ (or any descendant) is observed. Observing a common effect makes its
causes \emph{compete} — confirming one explains away the others.
\end{defbox}

The intuition is "explanatory work": if a wet lawn could be due to rain or sprinklers, observing wet
grass raises both. But once you confirm the sprinkler was on, it already accounts for the wet grass
— there is less explanatory work left for the rain to do, so rain's probability falls back. The two
causes compete for the same evidence.

\begin{exbox}[explaining away — the wet-grass numbers worked]
Priors $P(r{=}1)=0.20$, $P(s{=}1)=0.40$ (rain, sprinkler — marginally independent). CPT for wet
grass: $P(w{=}1\mid r,s)=0.95$ if rain \emph{or} sprinkler is on, $0.0125$ if neither. The joint of
each $w{=}1$ case is $P(r)P(s)P(w{=}1\mid r,s)$:
\[
(1,1)\!:0.076,\quad(0,1)\!:0.304,\quad(1,0)\!:0.114,\quad(0,0)\!:0.006
\;\;\Rightarrow\;\; P(w{=}1)=0.50.
\]
\kw{Observe wet grass}: $P(r{=}1\mid w)=\tfrac{0.076+0.114}{0.50}=0.38$ and
$P(s{=}1\mid w)=\tfrac{0.076+0.304}{0.50}=0.76$ — both rise above their priors.
\kw{Also observe $s{=}1$}: $P(r{=}1\mid w,s{=}1)=\tfrac{0.076}{0.076+0.304}=0.20$ — rain drops back
to its prior, ``explained away'' by the confirmed sprinkler.
\end{exbox}

d-separation is the calculus that lets you read \emph{any} conditional independence directly off the
graph by walking every path and applying one of three rules at each middle node.

\begin{defbox}[d-separation]
$a\perp b\mid z$ iff \kw{every} path between them is blocked.
\kw{Chain} $a\to z\to b$ — blocked iff $z$ observed.
\kw{Fork} $a\leftarrow z\to b$ — blocked iff $z$ observed.
\kw{Collider} $a\to z\leftarrow b$ — open iff $z$ \emph{or any descendant} observed
(\kw{reversed} from chain/fork).
\end{defbox}

For a chain, information about $a$ flows through $z$ to $b$; pinning $z$ down (observing it) stops
that flow. For a fork, $z$ is a common cause, so if you do not know $z$ then $a$ tells you about $z$
which tells you about $b$ — observing $z$ severs the link. The collider is the odd one out:
observing the middle node \emph{opens} the path (explaining-away). The standard memory aid is to
think of information as a ball bouncing — it bounces only off arrow-heads, so the V's two arrowheads
at $z$ only conduct when $z$ is observed. The "or any descendant" clause matters: if a collider's
child is observed, partial information about $z$ leaks through it and the path is still open. Two
RVs are d-separated iff \emph{every} path between them is blocked — if any single path is open they
are dependent.

\begin{notebox}[why the rules hold — idea only (proofs non-examinable)]
\kw{Chain}: $P(a,z,b)=P(a)P(z|a)P(b|z)$; divide by $P(z)$, rewrite $P(a)P(z|a)=P(z)P(a|z)$, $P(z)$
cancels $\Rightarrow P(a,b|z)=P(a|z)P(b|z)$. \kw{Collider} $z$ unobserved: $\sum_z P(z|a,b)=1$ so
$P(a,b)=P(a)P(b)$; $z$ observed: the term $P(z|a,b)$ binds $a$ and $b$ and cannot be split into a
factor in $a$ alone times a factor in $b$ alone — so they are dependent. The takeaway: marginal
independence and conditional independence are genuinely different regimes, and the graph encodes
both.
\end{notebox}

% =====================================================================
\section{Week 23 — Message Passing}

Week 22 gave us graphical models as compact representations of joint distributions; Week 23 supplies
the algorithmic engine that makes them usable. Once you have a model and some evidence, the question
you actually want answered is "what is the probability of this one RV?" — and answering it naively
means summing the joint over every assignment of every other variable, which is exponential.
Message passing is the dynamic-programming trick that brings this down to polynomial time, exactly
on chains and trees.

\begin{setupbox}[the inference problem]
Given a graphical model + evidence, compute the marginal / \kw{belief} of any RV. Naive enumeration
is \kw{$O(n^N)$} ($n$ states, $N$ RVs) — exponential. \kw{Message passing} (sum--product) does it in
\kw{$O(Nn^2)$} for chains and trees.
\end{setupbox}

\subsection{Markov chains}

A Markov chain is the simplest graphical model with temporal structure: a 1D chain of RVs in which
each new state depends only on the previous one (the Markov property — "the state has no memory of
the past"). It is cleanest to understand by contrast with the finite state machine it generalises.

\begin{defbox}[FSM vs Markov chain]
\kw{Finite state machine} — deterministic: an \kw{external input} (an observed outside event) chooses
the next state with certainty; no probabilities. \kw{Markov chain} — probabilistic: advances
autonomously, the next state drawn from the \kw{transition matrix} $T$ ($T_{ab}=P(s_a\to s_b)$, rows
are categorical distributions).
\end{defbox}

\begin{notebox}[an FSM ``runs on'' a Markov chain]
Replace the explicit external inputs with \kw{how often they occur} — the input frequencies become
the transition probabilities. The FSM topology (which transitions are allowed) is the skeleton; $T$
is the engine that now advances the system on its own. A Markov chain = an autonomous, input-free FSM.
\end{notebox}

Three operations on a Markov chain matter for the exam — learning the matrix, sampling from it, and
extending its memory.

\begin{itemize}
  \item \kw{Learning $T$}: each row of $T$ is a categorical distribution, so you initialise every
  cell to 1 (a uniform Dirichlet prior, equivalently $+1$ Laplace smoothing), count every observed
  transition $s_a\to s_b$, and normalise each row to sum to 1. The result is the MAP estimate;
  Dirichlet is the conjugate prior for the categorical, mirroring the Beta--Bernoulli conjugacy of
  Week 22.
  \item \kw{Sampling}: draw a start state from $P_{\text{start}}$ (fitted from training first-states),
  draw a length $L\sim\text{Poisson}(\bar L)$ from the mean training length, then step
  $x_i\sim T_{x_{i-1},\cdot}$. This is used for generative modelling and as a sanity check — bad
  samples reveal a bad model.
  \item \kw{$n$-gram}: a 1-character-memory chain produces phoneme-like nonsense ("ackilin,
  ckleryhtz") because each adjacent letter-pair looks plausible while the word as a whole drifts.
  Crucially the fix is a longer \emph{generation-time} memory window — make the state the last $n$
  characters — \kw{not} more training data, since at generation time the model has already forgotten
  everything but the last character. The cost is $26^n$ states, so each step up demands
  exponentially more data to fit reliably (the curse of dimensionality again).
\end{itemize}

\subsection{Inference and the message rule}

\kw{Marginal} $P(x_i)=\sum_{\text{others}}P(x_1,\dots,x_N)$; \kw{belief} = the same with evidence
applied. \kw{Notation} $b_j(x_j)$: subscript $j$ = \emph{which node}, argument $x_j$ = \emph{which
state} — it is a \emph{vector}, the global updated distribution for node $j$ (not ``predict the next
state'', which would be $P(x_j\mid x_{j-1})$).

Naive marginalisation is $O(n^N)$ — a nested sum over every assignment of the other RVs (for $n=26$,
$N=10$ that is $\approx 1.4\times10^{14}$ operations, hopeless). The \kw{fix} rests on one
observation: most terms in the joint do not depend on most summation variables, so push each $\sum$
inside the products that do not depend on its variable. On a chain
$b_3=\sum_{x_1,x_2,x_4,x_5}\prod P$, the sum over $x_1$ touches only $P(x_1)P(x_2|x_1)$, so
$\sum_{x_1}$ moves inward and collapses to a function of $x_2$ alone — once you have summed over
every value of a variable, the result simply no longer depends on it. Continue down the chain and
every sum becomes a single-variable sum with a two-term summand. This is \kw{dynamic programming}:
cache each partial sum (a \kw{message}) and reuse it. Each is an $O(n^2)$ sum, $O(N)$ of them
$\Rightarrow O(Nn^2)$ total — ten orders of magnitude cheaper for the example above.

\begin{defbox}[sum--product message passing]
A \kw{message} $m_{i\to j}(x_j)$ is a vector, one entry per state of $j$. Chain/tree rule:
\[
m_{i\to j}(x_j)=\sum_{x_i}U(x_i)\,F(x_i,x_j)\!\!\prod_{k\in N_i\setminus j}\!\! m_{k\to i}(x_i).
\]
\kw{Belief}: $b_j(x_j)=U(x_j)\prod_{i\in N_j}m_{i\to j}(x_j)$ — local unary $\times$ \emph{all}
incoming messages.
\end{defbox}

\begin{notebox}[reading the message rule — right to left]
$\prod m_{k\to i}$ = the \kw{inbox} (all messages into $i$ \emph{except} from $j$ — excluding $j$
stops an echo chamber); $U(x_i)$ = \kw{local context} ($i$'s own unary factor); $F(x_i,x_j)$ = the
\kw{bridge} translating $i$'s belief into $j$'s states; $\sum_{x_i}$ = \kw{sum out} $x_i$, leaving a
clean vector in $x_j$. Slogan: \kw{multiply incoming, sum out the rest}. Edges are undirected (a
forward + a backward pass cover both directions) but a single message never feeds back to its sender.
\end{notebox}

\begin{exbox}[message passing on a 5-node chain — computing $b_3$]
Chain $x_1\!-\!x_2\!-\!x_3\!-\!x_4\!-\!x_5$, query the middle node.
\kw{Forward}: $m_{1\to2}(x_2)=\sum_{x_1}U(x_1)F(x_1,x_2)$, then
$m_{2\to3}(x_3)=\sum_{x_2}U(x_2)F(x_2,x_3)\,m_{1\to2}(x_2)$.
\kw{Backward}: $m_{5\to4}(x_4)=\sum_{x_5}U(x_5)F(x_4,x_5)$, then
$m_{4\to3}(x_3)=\sum_{x_4}U(x_4)F(x_3,x_4)\,m_{5\to4}(x_4)$.
\kw{Belief}: $b_3(x_3)=U(x_3)\,m_{2\to3}(x_3)\,m_{4\to3}(x_3)$ — the node's own term times
left-context times right-context. One forward + one backward pass ($2(N{-}1)$ messages, each
$O(n^2)$) yields \emph{every} node's belief — $O(Nn^2)$ total, against $O(n^N)$ for naive
enumeration.
\end{exbox}

The reuse is the key economy. Once every message on the chain has been computed, the belief at
\emph{every} node falls out as a product of incoming messages — you do not redo the work per node.

\begin{itemize}
  \item \kw{Two messages per edge} (one each direction, because the two directions carry different
  information — left-side context vs right-side context). One forward pass plus one backward pass
  computes all $2(N{-}1)$ messages, after which every node's belief comes "for free" — $O(Nn^2)$
  total, versus $O(N\cdot n^N)$ for naive marginalisation at every node.
  \item \kw{Factor graph} (two message types, because circles and squares alternate):
  variable$\to$factor $m_{v\to F}=\prod_{G\ne F}m_{G\to v}$ (a variable forwards what every
  \emph{other} factor told it); factor$\to$variable
  $m_{F\to v}=\sum_{x\setminus v}F(x)\prod_{u\ne v}m_{u\to F}$ (a factor combines its own function
  with incoming variable messages, then marginalises out everything but the destination). The two
  rules are still "multiply then sum" — the algorithm's name, sum--product, comes from exactly this
  pair of operations.
\end{itemize}

\begin{notebox}[HMM — the voice-recognition example]
Hidden phonemes $x_i$, observed audio frames $y_i$. Two factors: \kw{unary} $U_i(x_i)\propto P(y_i|x_i)$
— the per-frame acoustic guess (instantaneous, no context); \kw{pairwise} $P_{i,i+1}$ — the phoneme
transition statistics (temporal context). Belief $b_j=U_j\times$ forward message $\times$ backward
message = local audio fit fused with left- and right-context. This model is a \kw{Hidden Markov Model}.
\end{notebox}

\begin{trapbox}[message passing is exact only on trees / chains]
The whole method relies on each message summarising "everything I know from my side." On a tree
every edge cleanly partitions the graph into two sides, so a message can summarise one side
perfectly. A loop breaks that partition — information flowing round the loop returns to its origin
and double-counts. Loopy graphs need approximate variants (loopy belief propagation — iterate and
hope for convergence; junction tree — cluster RVs to break loops, exact but exponential in cluster
size; sampling). Exact and $O(Nn^2)$ only for tree-structured factor graphs.
\end{trapbox}

\begin{center}\large\bfseries\color{defBd} PART VIII — Support Vector Machines (Weeks 25–26)\end{center}

% =====================================================================
\section{Week 25 — Linear Support Vector Machines}

For linearly separable data there are infinitely many hyperplanes that classify the training set
perfectly — and squared error or accuracy cannot tell them apart, just as in the Week 21 four-point
regression. Some of those separators sit right against one class and would misclassify the next
point that arrives; others keep a comfortable gap. The SVM operationalises "comfortable gap" as a
precise criterion: pick the separator whose distance to the nearest training point is largest. That
single idea — maximise the margin — turns classification into a clean convex optimisation problem.

\begin{setupbox}[hard-margin SVM]
Binary classification, labels $y_i\in\{-1,+1\}$. A \kw{separating hyperplane}
$H=\{x:w^\top x+b=0\}$ satisfies $y_i(w^\top x_i+b)>0$ for all $i$. Infinitely many exist — the SVM
picks the \kw{maximum-margin} one (furthest from the closest point of either class).
\end{setupbox}

The product $y_i(w^\top x_i+b)$ is positive exactly when the prediction sign matches the label: on
the $+1$ side the score is positive and $y_i=+1$, on the $-1$ side the score is negative and
$y_i=-1$, and a negative times a negative is positive. So "$>0$ for all $i$" is a compact way to
say "every point correctly classified."

\begin{pitfallbox}[SVM labels are $\{-1,+1\}$, not $\{0,1\}$]
The constraint $y_i(w^\top x_i+b)\ge1$ relies on $y_i$ to flip the sign. With the logistic-regression
convention $y_i=0$, the left side collapses to $0$ and $0\ge1$ is impossible — the constraint is
destroyed. Always convert labels first.
\end{pitfallbox}

\subsection{Canonical rescaling and the margin}

The hyperplane $w^\top x+b=0$ is \kw{scale-invariant} — $(cw,cb)$ for $c>0$ is the same boundary, so
the raw score $y_i(w^\top x_i+b)$ can be inflated freely. \kw{Canonical rescaling}: use that freedom
to divide $(w,b)$ by the closest point's score, so $y_i(w^\top x_i+b)\ge1$ for all $i$ with equality
on the closest points. The ``$1$'' is an artificial calibration of the scale, not a law of nature.

\begin{notebox}[$w^\top x_i+b$ can never be 0 for a training point]
Since $|y_i|=1$, the constraint $y_i(w^\top x_i+b)\ge1$ is equivalent to $|w^\top x_i+b|\ge1$ — the
score has magnitude $\ge1$, so no training point can sit on $H$. Pure algebra; the margin-band
geometry is its consequence.
\end{notebox}

With the scale fixed, the margin becomes a concrete distance you can derive (both proofs are past
paper Q6(a)/(b) and worth knowing cold).

\begin{itemize}
  \item \kw{$w\perp H$}: for any two points $x_1,x_2\in H$, both satisfy $w^\top x+b=0$; subtracting
  gives $w^\top(x_1-x_2)=0$. The vector $x_1-x_2$ lies \emph{within} the hyperplane, and its dot
  product with $w$ is zero — so $w$ is perpendicular to every direction in $H$, i.e.\ $w$ is the
  normal vector. This is why the shortest path between the two margin hyperplanes runs along $w$.
  \item \kw{Margin} $m=2/\|w\|$ — the distance between $H_+:w^\top x+b=1$ and $H_-:w^\top x+b=-1$.
  \emph{Derivation}: take $z_1\in H_+$, $z_2\in H_-$; since the shortest path is along the normal,
  write $z_1-z_2=\lambda w$. Substituting $z_1=z_2+\lambda w$ into $w^\top z_1+b=1$ and using
  $w^\top z_2+b=-1$ gives $-1+\lambda\|w\|^2=1\Rightarrow\lambda=2/\|w\|^2$, so
  $m=\|z_1-z_2\|=\lambda\|w\|=2/\|w\|$. Therefore maximising the margin $m$ is exactly the same as
  minimising $\|w\|$, and equivalently minimising $\tfrac12\|w\|^2$ — the half and the square are
  conventions that make the gradient clean.
  \item \kw{Functional margin} $y_i(w^\top x_i+b)$ = the raw signed score, positive iff correctly
  classified but \emph{scale-dependent} (inflating $(w,b)$ inflates it); \kw{geometric margin}
  $1/\|w\|$ = the true perpendicular distance from a margin hyperplane to $H$, \emph{scale-independent}.
  The optimisation cares about the geometric one; the functional one is just a calibration artefact.
\end{itemize}

\subsection{Optimisation, Lagrangian, KKT}

\begin{defbox}[hard-margin SVM problem]
\[
\min_{w,b}\;\tfrac12 w^\top w \quad\text{s.t.}\quad y_i(w^\top x_i+b)\ge1\ \ \forall i.
\]
Convex \kw{quadratic program} (quadratic objective + linear constraints) $\Rightarrow$ unique optimum.
\end{defbox}

\begin{notebox}[a linear SVM is the last layer of a deep network]
A full deep net is non-convex (stacked non-linear layers) $\Rightarrow$ needs SGD/Adam and
learning-rate tuning. Strip the hidden layers, keep only the final linear classifier, and you get the
SVM — a convex QP where every local minimum is global, so a QP solver suffices.
\end{notebox}

To solve a constrained problem like this you roll the constraints into the objective with a
\kw{Lagrangian}, one multiplier $\mu_i$ per constraint:
$L(w,b,\mu)=\tfrac12 w^\top w+\sum_i\mu_i(1-y_i(w^\top x_i+b))$, with $\mu_i\ge0$. (Here the
constraint is written in the standard $g_i\le0$ form $g_i=1-y_i(w^\top x_i+b)$.)

\begin{notebox}[why $\mu_i\ge0$ — three readings]
\kw{Penalty}: $\mu_i g_i$ must \emph{raise} $L$ when a constraint is violated ($g_i>0$) — a negative
$\mu_i$ would reward cheating. \kw{Normal force}: the constraint ``wall'' can only push, not pull.
\kw{Shadow price}: $\mu_i$ = how much the optimum worsens if the constraint is tightened — tightening
never helps, so $\mu_i\ge0$; an inactive constraint has price $0$.
\end{notebox}

\begin{defbox}[KKT conditions]
$\theta^\ast$ is optimal iff $\exists\,\mu^\ast\ge0$ with: \kw{stationarity} $\nabla_\theta L=0$;
\kw{primal feasibility} $g_i(\theta^\ast)\le0$; \kw{dual feasibility} $\mu_i^\ast\ge0$;
\kw{complementary slackness} $\mu_i^\ast\,g_i(\theta^\ast)=0$ (per constraint).
\end{defbox}

The four conditions are each necessary, and because the SVM problem is convex they are jointly
sufficient — a $(\theta^\ast,\mu^\ast)$ satisfying all four is genuinely optimal. Complementary
slackness is the load-bearing one: for each $i$, either the multiplier is zero or the constraint is
tight, never both strict.

Applying KKT to the SVM, \kw{stationarity in $w$} gives $w^\ast=\sum_i\mu_i^\ast y_i x_i$ — the
optimal hyperplane is a signed linear combination of the input points. \kw{Stationarity in $b$}: of
the Lagrangian, only the term $-b\sum_i\mu_i y_i$ contains $b$, so
$\partial L/\partial b=-\sum_i\mu_i y_i=0\Rightarrow\sum_i\mu_i^\ast y_i=0$ — the positive- and
negative-class multipliers must balance.

\subsection{Support vectors and sparsity}

\begin{trapbox}[why the SVM solution is sparse — the full exam answer]
Two KKT lines combine. \kw{Stationarity}: $w^\ast=\sum_i\mu_i^\ast y_i x_i$. \kw{Complementary
slackness}: any point strictly inside its half-plane has $g_i<0$, which forces $\mu_i^\ast=0$. Those
points \kw{drop out of the sum} — $w^\ast$ is built only from the boundary points
($y_i(w^{\ast\top}x_i+b^\ast)=1$), the \kw{support vectors}. Sparsity is not a bonus property; it is
the direct algebraic consequence of complementary slackness zeroing $\mu_i$ for every non-margin point.
\end{trapbox}

\begin{itemize}
  \item Only support vectors define the hyperplane: $w^\ast=\sum_{i\in\mathcal S}\mu_i^\ast y_i x_i$
  is a sum over the margin points alone, so removing any non-support point leaves the optimum
  unchanged. The bias is then recovered from any single support vector via
  $b^\ast=y_{i_0}-w^{\ast\top}x_{i_0}$, since that point sits exactly on its margin hyperplane.
  (In practice the equality $y_i(w^{\ast\top}x_i+b^\ast)=1$ is never met exactly in floating point,
  so support vectors are identified within a small tolerance.)
  \item \kw{Three payoffs of sparsity}: \kw{robustness} — the boundary is fixed by the hardest
  examples right at the margin, so adding or moving easy interior points cannot perturb it;
  \kw{test-time memory} — only the support vectors need storing, the rest of the training set can be
  discarded; and it is what makes the kernel trick (Wk26) efficient, since predictions reduce to a
  short sum over support vectors.
  \item \kw{Multi-class}: the binary SVM extends via one-vs-rest — train $K$ binary SVMs, the
  $k$-th treating class $k$ as positive and all others as negative, then predict the class whose
  score $w_k^{\ast\top}x+b_k^\ast$ is largest.
\end{itemize}

% =====================================================================
\section{Week 26 — Non-linear SVM}

The Week 25 hard-margin SVM assumed the data was linearly separable. Real data rarely is, in two
distinct ways: it may need a linear boundary but with a few points on the wrong side, or it may need
a genuinely curved boundary. Week 26 lifts both restrictions — slack variables for the first,
feature maps and kernels for the second — and in doing so rebuilds the SVM as the \emph{dual}
problem, which is what makes the kernel trick possible.

\begin{setupbox}[two extensions of the hard-margin SVM]
\kw{Soft-margin} — handle non-separable data with slack variables. \kw{Feature map / kernel} —
non-linear decision boundaries.
\end{setupbox}

\subsection{Soft-margin SVM and hinge loss}

If no hyperplane can satisfy $y_i(w^\top x_i+b)\ge1$ for every point, the hard-margin problem is
infeasible. The fix is to give each point a non-negative \kw{slack} $\xi_i$ that lets it violate the
constraint — but slack is not free, or it could be made arbitrarily large and the constraint would
become vacuous, so it is paid for in the objective.

\begin{defbox}[soft-margin SVM]
\[
\min_{w,b,\xi}\;\tfrac12 w^\top w + C\sum_i\xi_i \quad\text{s.t.}\quad
y_i(w^\top x_i+b)\ge1-\xi_i,\ \ \xi_i\ge0.
\]
\kw{Slack} $\xi_i\ge0$ allows constraint violation; $C$ = the error/margin trade-off hyperparameter.
\end{defbox}

Reading the slack: $\xi_i=0$ is hard-margin behaviour (correctly classified, outside the margin);
$0<\xi_i\le1$ means the point is correct but has crept inside the margin band; $\xi_i>1$ means it is
outright misclassified.

\begin{notebox}[$C$ vs $\lambda$ — a structural flip]
A standard regularised model is $\text{Loss}+\lambda\cdot\text{Reg}$ — the penalty is on the
\emph{regularisation} term. The soft-margin SVM is
$\underbrace{\tfrac12 w^\top w}_{\text{reg}}+C\underbrace{\sum\xi_i}_{\text{loss}}$ — the penalty $C$
is on the \emph{loss} term. Hence they are inverses: large $C$ punishes errors hard $\Rightarrow$
small margin, near-hard-margin (\emph{less} regularisation, small $\lambda$); small $C$ $\Rightarrow$
wide margin, more slack tolerated (more regularisation).
\end{notebox}

The slack variables can be eliminated entirely. For fixed $(w,b)$ the constraint forces
$\xi_i\ge1-y_i(w^\top x_i+b)$ when that right-hand side is positive, and the objective is increasing
in $\xi_i$, so the optimal choice is the smallest feasible value
$\xi_i=\max(0,1-y_i(w^\top x_i+b))$. Substituting it back gives the unconstrained form
$\min_{w,b}\tfrac12 w^\top w + C\sum_i\max(0,1-y_i(w^\top x_i+b))$. The \kw{hinge loss}
$L_{\text{hinge}}(y\hat y)=\max(0,1-y\hat y)$ is the \kw{per-point loss term only} — \emph{not} the
whole objective (which is regularisation $+\,C\times$ hinge). It is zero for $y\hat y\ge1$ (the
model is "done" with a confidently correct point) and rises linearly for $y\hat y<1$; note it
penalises points that are correctly classified but not confident enough — anything inside the
margin still incurs loss. The name is geometric: flat at zero, then bending linearly, like a hinge.

\begin{trapbox}[hinge vs 0--1 loss]
\kw{0--1 loss} ($1$ if $y\hat y<0$, else $0$) is the true classification metric, but it is
\kw{non-convex and non-differentiable} — a step function — so optimising it directly is NP-hard.
\kw{Hinge} is a \kw{convex upper bound} on 0--1: it lies above the step everywhere and the two
coincide at $y\hat y=0$ and for $y\hat y\ge1$. Minimising hinge therefore minimises an upper bound
on the misclassification rate, and being convex it is solvable by standard convex optimisation, with
a decision boundary close to the (inaccessible) 0--1 optimum. Other convex surrogates exist
(logistic, squared hinge); hinge is the canonical SVM choice.
\end{trapbox}

\subsection{Soft-margin KKT and feature maps}

The soft-margin Lagrangian carries a second multiplier $\lambda_i$ for the $\xi_i\ge0$ constraint.
Stationarity in $\xi_i$ gives $C-\mu_i-\lambda_i=0$, i.e.\ $\lambda_i=C-\mu_i$; combining with
$\lambda_i\ge0$ and $\mu_i\ge0$ yields the \kw{box constraint} $0\le\mu_i\le C$ — the one structural
change from hard-margin (which had only $\mu_i\ge0$). The upper bound caps how much any single
training point can influence the optimum. Reading off the three regimes per point:
\begin{itemize}
  \item \kw{$\mu_i=0$}: outside the margin, well-classified — \emph{not} a support vector, ignored;
  most points are here.
  \item \kw{$0<\mu_i<C$, $\xi_i=0$}: exactly on the margin — a standard support vector, and these
  are the points used to recover $b^\ast$.
  \item \kw{$\mu_i=C$, $\xi_i>0$}: inside the margin or misclassified — a "hard" support vector,
  whose slack absorbs the violation.
\end{itemize}
In every case the support vectors are the challenging examples; comfortable interior points
contribute nothing.

For curved boundaries the move is a \kw{feature map} $\phi:\mathbb R^d\to\mathbb R^D$ ($D\gg d$):
lift the inputs into a higher-dimensional space and run a \emph{linear} SVM there. A linear boundary
in feature space corresponds to a non-linear surface in input space — for instance a $\phi$ carrying
an $x_1x_2$ term separates the XOR pattern, since $x_1x_2$ is positive in two opposite quadrants and
negative in the other two, and the resulting input-space boundary is a hyperbola. The SVM algorithm
is unchanged; only the representation is. This hand-crafting of features is exactly the job that
deep learning later automates — which is why deep learning is also called representation learning.

\subsection{The dual and the kernel trick}

The \kw{dual} is obtained by minimising the Lagrangian over the primal variables and then maximising
what remains over the multipliers. The primal has $d+1+n$ variables; the dual has only $n$, and —
more importantly — it expresses everything through inner products. Tracing the collapse: distribute
$\mu_i$ through the Lagrangian, then apply the two stationarity results. The bias term
$-b\sum_i\mu_i y_i$ vanishes because $\sum_i\mu_i y_i=0$. The cross term
$-w^\top(\sum_i\mu_i y_i\phi(x_i))$ becomes $-w^\top w$ because the bracket \emph{is} $w$. With
$\tfrac12 w^\top w-w^\top w=-\tfrac12 w^\top w$, and expanding the surviving $w^\top w$ into a double
sum:
\[
\max_{\mu}\;\sum_i\mu_i - \tfrac12\sum_{i,j}\mu_i\mu_j y_i y_j\,\phi(x_i)^\top\phi(x_j),
\qquad \sum_i\mu_i y_i=0,\ \ 0\le\mu_i\le C.
\]
Every primal variable is gone. The dual is a maximisation because each $g(\mu)$ is a valid
\emph{lower bound} on the primal optimum, and maximising chases the tightest (highest) such bound —
like guessing a number with only "higher/lower" feedback. \kw{Strong duality} holds for the SVM
(it is convex with a feasible interior, Slater's condition), so primal optimum $=$ dual optimum and
solvers may use whichever is convenient — usually the dual, both for the smaller variable count and
because it accommodates kernels.

\begin{notebox}[the kernel trick]
The dual — and prediction $w^{\ast\top}\phi(x)=\sum_{i\in S}\mu_i^\ast y_i\,\phi(x_i)^\top\phi(x)$ —
uses $\phi$ \kw{only as inner products} $\phi(x_i)^\top\phi(x_j)$, never $\phi$ on its own. So
replace each inner product by a \kw{kernel} $K(x_i,x_j)=\phi(x_i)^\top\phi(x_j)$ computed directly
from the raw inputs — $\phi$ is never formed. The RBF kernel $\exp(-\|x-z\|^2/2\tau^2)$ has an
\emph{infinite}-dimensional $\phi$ yet costs $O(d)$ per evaluation.
\end{notebox}

Prediction is kernelised too: substituting $w^\ast=\sum_{i\in S}\mu_i^\ast y_i\phi(x_i)$ makes the
test point interact with the support vectors only through inner products, and $b^\ast$ likewise
reduces to kernel evaluations — so the entire model, training and test, runs without ever forming
$\phi$.

\kw{The kernel trick — three advantages}: \kw{statistical} — non-linear boundaries lift accuracy on
data with non-linear structure; \kw{computational} — kernels avoid ever forming the (possibly
infinite-dimensional) feature space, an $O(d)$ kernel evaluation standing in for an inner product in
a space too large to construct; \kw{sparsity preserved} — the optimum still depends only on the
support vectors, so the model stays sparse in training points. Without the dual restructuring SVMs
would remain linear-only; the dual is what \emph{enables} the kernel trick.

\begin{center}\large\bfseries\color{defBd} PART IX — Kernel Methods (Week 29)\end{center}

% =====================================================================
\section{Week 29 — Kernel Methods}

\begin{setupbox}[the kernel idea]
Week 26 showed that the SVM dual — and its predictions — depend on the feature map $\phi$ \emph{only}
through inner products $\phi(x_i)^\top\phi(x_j)$. A \kw{kernel} is a function that returns exactly
that inner product, computed straight from the raw inputs, so the high-dimensional $\phi$ never has
to be built. Week 29 makes this rigorous: what \emph{is} a kernel, and how do you prove a given
function is one.
\end{setupbox}

\subsection{Why kernels exist — the cost they remove}

This is the ``why'' of the whole topic. A non-linear feature map deliberately \emph{inflates} the
dimensionality — taking, say, $10^3$ raw features up to $10^5$ or more — so that a problem which was
not linearly separable becomes separable by a hyperplane in the enlarged space. The catch: the SVM
dual sums $\phi(x_i)^\top\phi(x_j)$ over \emph{every pair} of training points, so the cost is
\kw{quadratic in the number of points $n$} and \kw{linear in the (exploded) feature dimension}.
Explicitly forming $\phi$ is therefore the bottleneck. \kw{A kernel removes exactly this cost}: it
computes the same inner product directly, without ever constructing $\phi$. Historically $n$ was
small (pre-big-data, pre-GPU), so the high feature dimension was the binding constraint — precisely
what kernels relieve.

\subsection{Formal definition}

\begin{defbox}[kernel]
A function $k:\mathcal X\times\mathcal X\to\mathbb R$ is a \kw{kernel} iff there exists a feature
space $\mathcal H$ (a \emph{Hilbert space} — a vector space, possibly infinite-dimensional, with a
well-defined inner product) and a feature map $\phi:\mathcal X\to\mathcal H$ such that
\[
k(x,x')=\langle\phi(x),\phi(x')\rangle_{\mathcal H}\qquad\text{for all }x,x'.
\]
\end{defbox}

\kw{Why this is powerful}: $\mathcal X$ \emph{need not be a vector space}. It can be a set of
documents, graphs, strings, or time series — objects with no natural geometry and no notion of
addition or dot product. The kernel \emph{imposes} an inner-product (similarity) structure on them,
which is what lets an SVM be run on non-vector data at all. The Hilbert-space requirement on
$\mathcal H$ is just "a vector space, possibly infinite-dimensional, where inner products make
sense"; when $\mathcal H$ is infinite-dimensional the inner product is an infinite sum, and the
Hilbert condition is precisely what guarantees that sum converges. The feature map is also \kw{not
unique}: e.g.\ the linear kernel admits both $\phi(x)=x$ in $\mathbb R$ and
$\phi(x)=(x,x)/\sqrt2$ in $\mathbb R^2$ — so "provide a feature map" always has several valid
answers, and two SVMs with the same kernel but different $\phi$ are computationally identical.

\subsection{Construction rules — building kernels from kernels}

Constructing a feature map by hand is hard for anything beyond the simplest kernels, so you rarely
do it. Instead you treat known kernels as building blocks and combine them with a small calculus of
rules — each rule is proved once by exhibiting the new feature map, after which it can be applied
freely:
\begin{itemize}
  \item \kw{Positive scalar}: $\alpha k$ is a kernel for $\alpha>0$, with new map
  $\sqrt\alpha\,\phi$. A \emph{negative} $\alpha$ fails because $\sqrt\alpha$ is not real and
  $\alpha\|\phi(x)\|^2<0$ would break the positive-definiteness any inner product must satisfy.
  \item \kw{Conic sum}: $\sum_j\alpha_j k_j$ is a kernel for $\alpha_j>0$ — concatenate (direct-sum)
  the feature maps, so the new map stacks the individual feature vectors.
  \item \kw{Product}: $k_1\cdot k_2$ is a kernel (stated without proof — it needs tensor-product
  Hilbert spaces — but usable as given).
  \item \kw{Mapping wrap}: $f(x)\,k(x,x')\,f(x')$ is a kernel for \emph{any} function $f$, with new
  map $f(x)\phi(x)$ in the same Hilbert space. This is the rule that handles the Gaussian kernel.
  \item \kw{Difference} $k_1-k_2$ is \emph{not} generally a kernel — and the way to see why is the
  next box.
\end{itemize}

\begin{trapbox}[the $x=x'$ trick — how to disprove a candidate kernel]
A genuine kernel must satisfy $k(x,x)=\langle\phi(x),\phi(x)\rangle=\|\phi(x)\|^2\ge0$ (the
positive-definiteness of any inner product). So if you can find a single $x$ with $k(x,x)<0$, no
feature map can exist and the candidate is not a kernel. This is how $k_1-k_2$ fails, and how
$\cos(x+x')$ is disproved.
\end{trapbox}

\subsection{Canonical kernels and their proofs}

The standard kernels are all proved the same way — rewrite them as a combination of simpler kernels
and invoke the construction rules — and each proof doubles as a description of the feature space it
implies.

\begin{itemize}
  \item \kw{Linear} $k=\langle x,x'\rangle$ — trivially a kernel ($\phi=$ identity).
  \item \kw{Polynomial} $k=(\langle x,x'\rangle+c)^m$, $c\ge0$. \emph{Proof}: binomial-expand into
  $\sum_i\binom mi c^{m-i}\langle x,x'\rangle^i$ — a conic sum (positive coefficients) of products of
  the linear kernel. Its feature space is all monomial feature interactions up to degree $m$.
  \item \kw{Cosine} $k=\cos(x-x')=\cos x\cos x'+\sin x\sin x'$ $\Rightarrow$ explicit map
  $\phi(x)=(\cos x,\sin x)$. \kw{Trap}: $\cos(x{+}x')$ is \emph{not} a kernel — at $x=x'=\pi/3$ it is
  $\cos(2\pi/3)=-0.5<0$.
  \item \kw{Exponential} $k=\exp(\langle x,x'\rangle)$. \emph{Proof}: Taylor-expand
  $e^z=\sum z^i/i!$ — a conic sum of polynomial kernels. Feature space is infinite-dimensional.
  \item \kw{Gaussian / RBF} $k=\exp(-\|x-x'\|^2/2\sigma^2)$. \emph{Proof}: expand
  $\|x-x'\|^2=\|x\|^2+\|x'\|^2-2\langle x,x'\rangle$, giving
  $k=f(x)\cdot\exp(\langle x,x'\rangle/\sigma^2)\cdot f(x')$ with $f(x)=\exp(-\|x\|^2/2\sigma^2)$ —
  a mapping wrap of the exponential kernel. Its feature space is \kw{infinite-dimensional}: $\phi$
  can never be formed, yet $k$ is a one-line computation. \emph{That} is the magic of kernels.
\end{itemize}

\begin{notebox}[the RBF bandwidth $\sigma$ — a tunable knob, and its limitation]
\kw{Small $\sigma$}: each training point's influence is highly localised $\Rightarrow$ a wiggly
decision boundary that overfits. \kw{Large $\sigma$}: broad influence $\Rightarrow$ a smooth,
near-linear boundary that underfits. $\sigma$ is a hyperparameter, cross-validated like $C$ — and a
poor choice undermines the model regardless of how much data you have.
\end{notebox}

\subsection{Kernel methods vs deep learning — strengths and limitations}

The cleanest way to place kernel methods is against the technology that displaced them. Both produce
non-linear classifiers, but they get there in opposite ways, and the trade-off explains why each
dominated its era.

\begin{center}
\renewcommand{\arraystretch}{1.2}
\begin{tabularx}{\linewidth}{l X X}
\toprule
& \thead{Kernel SVM} & \thead{Deep neural network} \\
\midrule
Feature representation & fixed by the kernel choice & learned end-to-end from data \\
Optimisation & convex (QP) — one global optimum & non-convex — many local optima \\
Sample efficiency & high — works on small datasets & low — needs large labelled datasets \\
Scaling in $n$ & $O(n^2)$ or worse & $O(n)$ per pass \\
Non-linearity from & the kernel feature map & non-linear activation functions \\
\bottomrule
\end{tabularx}
\end{center}

\kw{Strengths of kernel methods}: convex (a unique optimum, no initialisation sensitivity);
sample-efficient; and they handle high-dimensional \emph{and} non-linear data \emph{simultaneously}
(past-paper MCQ — the answer is ``both''). \kw{Limitations}: the feature map is \emph{fixed} — chosen
by the practitioner, not learned, so it cannot adapt to the specific problem; and training scales
\kw{$O(n^2)$ or worse} in the number of points, so kernels do not reach massive datasets. Deep
learning's central advance is to \emph{learn} the feature representation — paying for it with a
non-convex problem and a large data requirement. Kernels dominated ML from the early 1990s to the
mid-2010s; deep learning since.

% =====================================================================
\section*{One-line recall checklist}

\begin{tldrbox}[Keyword prompts]
\begin{itemize}[leftmargin=*]
  \item Tree training: greedy per node; global-optimal tree NP-hard.
  \item IG \& Gini: same weighted-child template; IG $\to$ variance reduction for regression, Gini cannot.
  \item 0--1 loss useless inside a tree (blind to ``less wrong'') and for GD (zero gradient, NP-hard).
  \item Quantised inputs split at bin boundaries; categorical = one-left-rest-right, $O(n)$.
  \item Overfit = high train / low test acc.
  \item Naive variance biased $\to$ divide by $n-1$.
  \item MSE = Bias$^2$ + Variance + irreducible $\sigma^2$; proof non-examinable, components examinable.
  \item Regularisation $+\lambda\Omega(w)$: shrinks weights $\to$ less variance, slight bias; L1 = sparsity.
  \item Three-way split train/validation/test; never tune on test.
  \item Bootstrap: size $n$, $S$ datasets, $\approx1/e$ OOB; random subspace $\sqrt n$ / $n/3$;
  bagging variance $V/S$.
  \item Bayes classifier uses true $P$; noise floor = Bayes error = aleatoric uncertainty.
  \item Bayes-optimal threshold $x_0$ = curve crossover; errors are region integrals; $\hat x$ adds excess risk.
  \item Convex surrogates: hinge, logistic; squared loss over-penalises confident-correct.
  \item Confusion matrix rows = predicted, cols = actual; balanced accuracy = \emph{mean} of recalls;
  accuracy misleads on imbalanced data.
  \item Generic optimiser: direction $p=-\nabla f$, step $\alpha$, termination ($T_{\max}$ / $\|\nabla f\|$ / progress).
  \item Stationary point $\ne$ minimum unless convex; convexity is about $L(w)$, not the data.
  \item SGD update $w_t-2\alpha(w^\top x_{i_s}-y_{i_s})x_{i_s}$; mini-batch averages $B$ gradients.
  \item BGD: convex + step size $\Rightarrow$ global; SGD only with decaying step.
  \item Hessian PSD $\iff$ convex; PD $\Rightarrow$ unique min; $(X^\top X)^{-1}$ needs full column rank.
  \item Near-collinear $\to$ $\det\to0$ $\to$ ill-conditioned $\to$ amplified rounding noise.
  \item Anisotropic loss $\to$ GD zig-zags (still convex).
  \item Newton: quadratic loss $\to$ one shot; higher-order $\to$ iterate (quadratic convergence).
  \item Newton fails: singular Hessian / non-differentiable / far start; damp with $(H+\lambda I)$, $\lambda>-e_{\min}$.
  \item GD beats Newton at scale: $O(M)$ vs $O(M^2)$ memory, $O(M^3)$ inversion.
  \item Likelihood = P(observed data $\mid$ $w$), a function of $w$ — not P($w$ correct).
  \item ML + i.i.d.\ Gaussian noise $\equiv$ ordinary least squares.
  \item Posterior $\propto$ likelihood $\times$ prior; MAP maximises the posterior.
  \item MAP + Gaussian prior $\equiv$ ridge, regularisation strength $=\sigma^2$.
  \item $\sigma^2$ cancels in ML (no prior to balance), is the trust factor in MAP's tug-of-war.
  \item ML/MAP = point estimates, no uncertainty; full Bayesian = predictive distribution.
  \item Logistic regression = $\sigma(w^\top x)$, output is P($y{=}1$); sigmoid squashes $\mathbb R\to(0,1)$.
  \item $\sigma'=\sigma(1-\sigma)$; logit = log-odds, the inverse; score $z=w^\top x$ in log-odds units.
  \item Log-likelihood concave; gradient $X^\top(y-\sigma)$; no closed form, train iteratively.
  \item Decision boundary $w^\top x=0$, linear; threshold $0.5$ is a choice (FP vs FN trade-off).
  \item Separable data $\Rightarrow\|w\|\to\infty$, no finite maximiser $\Rightarrow$ regularise (= MAP Gaussian prior).
  \item Sigmoid = binary; softmax = multi-class generalisation.
  \item Scaling/outliers hurt linear, logistic, Bayesian regression; trees/forests are immune.
  \item Unsupervised = no labels; cluster = similar within, dissimilar across.
  \item K-means: init / assign / update / converge; minimises WCSS $J$; $O(KNM)$ per iteration.
  \item K-means converges (local min): every step lowers $J$, $J\ge0$, finite assignments.
  \item $L_2$ pairs with K-means (mean minimises it); $L_1$ $\to$ median = K-medians.
  \item K-means weaknesses: need $K$, random init, outliers, varying size/density, high $M$; fix init with best-of-$N$.
  \item Hierarchical = dendrogram; agglomerative bottom-up, divisive top-down; single linkage chains.
  \item Good clustering = low intra + high inter; Silhouette $(b{-}a)/\max(a,b)$ higher better, DBI lower better.
  \item Elbow method: WCSS-vs-$K$ bend picks $K$.
  \item Parametric = fixed \#params; non-parametric = \#params grows with $N$ (misleading name).
  \item Non-parametric: flexible + overfit-prone (TRUE); not fast, not small-data (FALSE).
  \item KNN = lazy learner, vote of $k$ nearest; $k$ is a hyperparameter; KNN $\ne$ K-means.
  \item KDE $\bar f_h(x)=\frac{1}{Nh}\sum G((x-x_i)/h)$; $h$ = bandwidth hyperparameter (Silverman / CV).
  \item Curse of dimensionality: samples needed scale $k^M$.
  \item Laplace = Gaussian at the mode; bad on multimodal / skewed / heavy-tailed targets.
  \item VI minimises KL divergence to the target; not Gaussian-restricted.
  \item Regularised loss = data fit + $\lambda R(\theta)$; $\lambda$ a validation-tuned hyperparameter.
  \item Four reasons to regularise: overfitting, ill-posed, interpretability, easier optimisation.
  \item Occam's razor = prefer the simplest model; regularisation is it written as an equation.
  \item Ridge $L_2$: closed form, shrinks but never zeros (not sparse). Lasso $L_1$: no closed form, sparse.
  \item $L_1$ sparse — diamond corners on axes + constant-$\lambda$ slope; $L_2$ — circle + fading $2w$ slope.
  \item Increasing $\lambda$ in lasso $\to$ sparser (TRUE); $N\uparrow\Rightarrow\lambda\downarrow$.
  \item Elastic net = $L_1{+}L_2$ blend; ridge = Gaussian prior, lasso = Laplace prior.
  \item Graphical model: nodes = RVs, edges = conditional (in)dependence; structure is an assumption.
  \item BN = directed $\prod P(\text{node}|\text{parents})$; factor graph = circles+squares; Markov net = undirected pairwise+unary.
  \item BN/MN $\to$ factor graph always; reverse not (pairwise limit + loss of explicitness).
  \item Shaded = observed; modes re-shade; latent = hidden vs hypothetical.
  \item Naive Bayes: $P(x|y)=\prod P(x_i|y)$, generative; MLE = count; predict in log space.
  \item Laplace smoothing = posterior \emph{mean} of a uniform Beta prior (not the mode).
  \item Collider/explaining away: marginally independent parents $\to$ conditionally dependent given the child.
  \item d-separation: chain/fork blocked iff middle observed; collider open iff middle/descendant observed (reversed).
  \item Message passing = inference for graphical models; naive $O(n^N)\to O(Nn^2)$ (dynamic programming).
  \item FSM = deterministic + external input; Markov chain = probabilistic; FSM runs on a MC by baking input frequencies into $T$.
  \item Learn $T$: count transitions + Laplace/Dirichlet, normalise rows; $n$-gram = longer memory, $26^n$ states.
  \item Belief $b_j(x_j)$: subscript = node, argument = state; global distribution, not next-state prediction.
  \item Message rule: multiply incoming (except sender) $\times$ unary $\times$ bridge factor, sum out; belief = unary $\times$ all incoming.
  \item Two messages per edge; one forward + one backward pass gives every node's belief.
  \item HMM = hidden states + unary acoustic factor + pairwise transition factor.
  \item Message passing exact only on trees/chains; loops $\to$ approximate.
  \item SVM labels $\{-1,+1\}$; constraint $y(w^\top x+b)\ge1$; $\{0,1\}$ breaks it.
  \item Hyperplane scale-invariant; canonical rescaling pins the closest point to 1; $|w^\top x+b|\ge1$.
  \item $w\perp H$; margin $=2/\|w\|$; maximise margin $\equiv$ minimise $\tfrac12\|w\|^2$.
  \item Functional margin scale-dependent; geometric margin $1/\|w\|$ scale-independent.
  \item SVM = convex QP = the last linear layer of a deep net.
  \item $\mu_i\ge0$: penalty / normal force / shadow price. KKT: stationarity, primal \& dual feasibility, complementary slackness.
  \item $w^\ast=\sum\mu_i y_i x_i$, $\sum\mu_i y_i=0$; complementary slackness zeros non-margin $\mu_i$ $\Rightarrow$ sparse, support vectors only.
  \item Soft-margin SVM: slack $\xi_i$, objective $\tfrac12 w^\top w+C\sum\xi_i$; $C$ on the loss term $\Rightarrow$ inverse of $\lambda$.
  \item Hinge loss $\max(0,1-y\hat y)$ = the per-point loss term only; convex upper bound on 0--1 loss.
  \item Soft-margin KKT: box constraint $0\le\mu_i\le C$; three regimes ($\mu{=}0$ / $0{<}\mu{<}C$ / $\mu{=}C$).
  \item Feature map $\phi$: linear SVM in feature space = non-linear boundary in input space.
  \item Dual uses $\phi$ only as inner products $\Rightarrow$ kernel trick: $K(x_i,x_j)=\phi(x_i)^\top\phi(x_j)$; RBF = infinite-dim $\phi$, $O(d)$ to evaluate.
  \item Kernel = $k(x,x')=\langle\phi(x),\phi(x')\rangle$ for some Hilbert space; works because the dual needs only inner products.
  \item $\mathcal X$ need not be a vector space (graphs, strings, documents); feature map non-unique.
  \item Construction rules: positive scalar, conic sum, product, mapping wrap; difference is NOT a kernel.
  \item Disprove a kernel with the $x{=}x'$ trick ($k(x,x)\ge0$ required); $\cos(x{-}x')$ yes, $\cos(x{+}x')$ no.
  \item Kernels: linear, polynomial (binomial), exponential (Taylor), Gaussian/RBF (mapping wrap, infinite-dim $\phi$).
  \item Kernel SVM = fixed feature map, convex, $O(n^2)$; deep net = learned features, non-convex, big data.
\end{itemize}
\end{tldrbox}

\end{document}
